% Which rotations join two Moser lattices into a 5-chromatic graph?
%
% Set at 10pt, which is what the target journals typeset at; at 11pt the same
% text runs to 27 pages and carries almost twice the over- and underfull boxes.
% Compiles with pdflatex, TeX Live 2023 or 2025.  No BibTeX: the bibliography is
% a manual thebibliography environment, which is what arXiv prefers for a short
% paper (no .bbl to mismatch).  No custom style files and no packages outside the
% TeX Live base.  Three figures are used -- spindle, law, orbits -- each a PDF in
% figures/ regenerated from the lattice itself by the script of the same name.
% figures/ holds five more from the 46-page draft that this version does not
% include; figures/README.md says which is which and how to rebuild them.

% Geombinatorics sets 11pt Times on a 5.5 x 8.5 inch page, no page numbers.
% That is a different document from the 10pt amsart above, so paper/geombinatorics.tex
% is a two-line wrapper that defines \geombinatorics and inputs this file.  There is
% deliberately no second copy of the manuscript: this project has paid three times
% for a number kept by hand in two places, and a whole second .tex is that defect
% with 1400 lines instead of one.
\ifdefined\geombinatorics
  \documentclass[11pt,reqno]{amsart}
\else
  \documentclass[10pt,reqno]{amsart}
\fi

\usepackage[T1]{fontenc}
\ifdefined\geombinatorics
  % geombina.uccs.edu/index.php/submissions-formatting, read rather than guessed.
  % "5.5 by 8.5 paper in 'Landscape' orientation" OR "the right side of a regular
  % 11 by 8.5 paper in landscape".  The second disambiguates the first: the right
  % half of an 11-wide, 8.5-tall sheet is 5.5 wide and 8.5 tall.  Portrait.
  % "Landscape" there names the sheet you print on, not the text block.
  % Margins are theirs verbatim: top, left, right 0.35"; bottom 0.5".
  \usepackage[paperwidth=5.5in,paperheight=8.5in,
              top=0.35in,left=0.35in,right=0.35in,bottom=0.5in]{geometry}
  % "Do NOT insert page numbers -- we will do it for you."
  \pagestyle{empty}
\fi
\usepackage[utf8]{inputenc}
\usepackage{amsmath}
\usepackage{amssymb}
\ifdefined\geombinatorics
  % Times, metrically New Roman.  newtx must come AFTER amssymb: loaded before,
  % both define \Bbbk and amssymb stops the run.
  \usepackage{newtxtext}
  \usepackage{newtxmath}
\fi
\usepackage{amsthm}
\usepackage{graphicx}
% The Geombinatorics set is the same three figures enlarged so their type is
% 11pt, as their checklist requires; figures/README.md says how to rebuild it.
% Same scripts, different output directory -- not a second copy of anything.
% law.pdf is the one figure included with a width key.  At 10pt that is right:
% it is drawn 342pt wide and the text block is 360pt, so the key fills the
% measure.  For Geombinatorics it would be a 1.4%% enlargement of a figure whose
% type is already exactly 11pt -- small, but it is scaling, and their
% illustrations are reproduced without reduction.  So: natural size there.
% NOT \includegraphics[\lawopt]{...}: keyval does not expand a macro in the
% option list, and \textwidth arrives there unexpandable.  The whole call is
% switched instead.
\ifdefined\geombinatorics
  \newcommand{\lawfig}{\includegraphics{law.pdf}}
\else
  \newcommand{\lawfig}{\includegraphics[width=\textwidth]{law.pdf}}
\fi
\ifdefined\geombinatorics
  \graphicspath{{figures/geom/}}
\else
  \graphicspath{{figures/}}
\fi
\usepackage[hidelinks]{hyperref}
% Two references are blog comments whose URLs are long and nearly unbreakable;
% justified without this, a short line stretches until its words stand alone.
\usepackage{xurl}
% A 40-character commit hash and the \texttt script names cannot break, so five
% lines ran into the right margin, the worst by 76pt.  Extra emergency stretch
% lets TeX loosen those lines instead of overflowing them; 2em is the least that
% clears all five.
\emergencystretch=2em
% \quad between the four basis vectors overflows the 4.8in Geombinatorics
% measure by 8pt.  A thinner separator there fixes it; at 10pt it stays \quad,
% so the deposited PDF's line breaking is untouched.
\ifdefined\geombinatorics
  \newcommand{\dsep}{\;}
\else
  \newcommand{\dsep}{\quad}
\fi
% A comma inside one $...$ is not a breakpoint, so a long inline list is a single
% unbreakable box -- 68pt over the 4.8in measure for the list of obstructed primes.
% \mc closes the maths, sets the comma in text where TeX may break, and reopens;
% at 10pt it stays an ordinary maths comma, so the deposited PDF is untouched.
% The list itself is written once either way.
\ifdefined\geombinatorics
  \newcommand{\mc}{$, $}
\else
  \newcommand{\mc}{,}
\fi
\newcommand{\doi}[1]{\href{https://doi.org/#1}{\texttt{#1}}}

\theoremstyle{plain}
\newtheorem{theorem}{Theorem}
\newtheorem{lemma}[theorem]{Lemma}
\newtheorem{proposition}[theorem]{Proposition}
\newtheorem{corollary}[theorem]{Corollary}
\newtheorem{conjecture}[theorem]{Conjecture}

\theoremstyle{definition}
\newtheorem{remark}[theorem]{Remark}

\newcommand{\Z}{\mathbb{Z}}
\newcommand{\Q}{\mathbb{Q}}
\newcommand{\R}{\mathbb{R}}
\newcommand{\C}{\mathbb{C}}
\newcommand{\F}{\mathbb{F}}
\newcommand{\Lat}{R}
\newcommand{\Ok}{\mathcal{O}}
\newcommand{\Gal}{\operatorname{Gal}}
\newcommand{\Lam}{\Lambda}
\DeclareMathOperator{\Tr}{Tr}
\DeclareMathOperator{\Nm}{N}
\DeclareMathOperator{\rank}{rank}

\begin{document}

\ifdefined\geombinatorics
% "Begin with the title in capital bold 14 pt, followed on the next line by the
% author(s)' name(s) in Non-Capital Bold 11 pt, followed by the address and
% e-mail non-bold 11 pt."  amsart's own title block is none of these, so it is
% replaced here rather than in the body: the 10pt build must not notice.
\makeatletter
\renewcommand{\maketitle}{%
  \thispagestyle{empty}%
  \begingroup
    \raggedright
    \fontsize{14}{17}\selectfont\bfseries
    A DIVISIBILITY OBSTRUCTION FOR JOINING TWO MOSER LATTICES\par
    \normalsize\bfseries Andrei-C\u{a}lin Nedelcu\par
    \normalfont\normalsize Bucharest, Romania\par
    \texttt{calin.nedelcu08@gmail.com}\par
  \endgroup
  \vspace{1.2\baselineskip}%
  % amsart keeps the abstract in \abstractbox, not in a macro; \@setabstract is
  % its own way of emptying that box, so the abstract survives the replacement.
  \@setabstract
  \vspace{1\baselineskip}%
}
\makeatother
\fi

\title[A divisibility obstruction for joining two Moser lattices]
      {A divisibility obstruction for joining\\ two Moser lattices}

\author{Andrei-C\u{a}lin Nedelcu}
\address{Bucharest, Romania}
\email{calin.nedelcu08@gmail.com}

\subjclass[2020]{05C15, 52C10, 05C62, 11R11}
\keywords{Hadwiger--Nelson problem, unit-distance graph, chromatic number of the
plane, Moser lattice, real quadratic field}

% Fixed, not \today: the PDF deposited at 10.5281/zenodo.22895457 is dated
% 21 September 2026, and with \today a rebuild could never reproduce it.
\date{September 21, 2026}

\begin{abstract}
Write $\Lat$ for the Moser lattice and $\theta_a$ for the rotation with
$\cos\theta_a = (2a-1)/(2a)$, rational $a > \tfrac14$, so that $a$ is the squared
radius at which
$\theta_a$ pairs a point with its own image. We prove a divisibility obstruction:
if the joint lattice $\Lat + \theta_a(\Lat)$ has rank eight and $v_2(a) < 2$---for
integer $a$, $4 \nmid a$---it
carries a $4$-colouring by a group homomorphism, so no subset of it is
5-chromatic at any size, and no 5-chromatic union $L \cup \theta_a(S)$ exists for
any $L, S \subseteq \Lat$. The restriction to $\Lat$ is real and not a
convenience: $\Lat$ is one member of a nested family of lattices
(Section~\ref{sec:notclaimed}), and Heule's $553$-vertex part already lives in
the next one up, so a join built there is outside everything below. The obstruction is decidable in closed form at every
rational $a$, before any graph is built; it kills every \emph{rank-eight} $a$
with $4 \nmid a$ and is silent otherwise---the rank-four exceptions of
Proposition~\ref{prop:rank} are untouched, as is every $a$ with $4 \mid a$.

Its engine is a structural fact: $3\Lat$ is the maximal order of
$\Q(\sqrt{33}, i\sqrt3)$. From that, $\Lat$'s two integral geometric
$4$-colourings are the reductions modulo the two primes above $2$, and occupancy
of the shell at radius $\sqrt a$ becomes a relative-norm condition, which the
arithmetic of $\Ok_V$ solves in closed form: $\Lat$ has a point at squared radius $a = n/b$ in
lowest terms if and only if $b \in \{1,3,9\}$ and $v_p(n)$ is even at every prime
that splits in $\Q(\sqrt{33})$ and stays inert in $\Q(\sqrt{-3})$. The rotations
surviving both conditions have natural density zero among the integers.

Against that, constructions. Chaining a monochromatic-pair gadget of Parts' gives
a $5$-chromatic join at $a = 64n^2/9$ for every $n$, an infinite family that
includes the one integer spindle our filters cannot decide. Seven spindles are
verified to join, three of them---$\theta_{16}$, $\theta_{28}$,
$\theta_{36}$---by constructions for which we have found no precedent; each of
the seven ships a DRAT certificate, checked, so that the unions are not
$4$-colourable is a proof and not a solver's word. A periodic
refinement kills eleven further integer candidates below $100$. The interface any such
construction must contain is treated in a companion paper and is not claimed
here. The record of $509$ is untouched.
\end{abstract}

\maketitle

\section{Introduction}\label{sec:intro}

The Hadwiger--Nelson problem asks for $\chi(\R^2)$, the least number of colours
needed so that no two points at distance $1$ share one. De~Grey \cite{deGrey}
proved $\chi(\R^2) \ge 5$ in 2018 with a 1581-vertex unit-distance graph;
Polymath16, Heule \cite{Heule} and Parts \cite{Parts} drove the vertex count down
to $509$, where it has stood since 2020.

The record is built out of one lattice, and so is everything here. Fix the
\emph{Moser lattice} $\Lat$, the subgroup of $(\C,+)$ generated by the $30$ unit
vectors
\[
  \exp\!\big(i(k\cdot 60^\circ + j\,\theta_3/2)\big),
  \qquad 0 \le k < 6, \quad \lvert j\rvert \le 2,
  \qquad \theta_3 = \arccos\tfrac56,
\]
the half-angle $\varphi = e^{i\theta_3/2} = \tfrac{\sqrt{33}}{6} +
\tfrac{\sqrt3}{6}i$ being needed for the twelve with $j$ odd.

\emph{Two groups carry this name, and this paper is about the larger one.} The
recent literature---\cite{Engel}, \cite{Ruhland}, \cite{Ducz},
\cite{Matolcsi}, \cite{DuczVarga}---writes
$M = \Z\langle 1, \omega_1, \omega_3, \omega_1\omega_3 \rangle$, of rank $4$
with $18$ unit vectors, and calls that the Moser lattice. The group the record
is built in is the larger $\Lat$ above: it contains $M$ with index $9$ and has
$30$ unit vectors. What separates them is conjugation. $M$ is not stable under
it---$\omega_3 + \bar\omega_3 = 5/3 \notin M$, so $\bar\omega_3 \notin M$---while
$\Lat$ is, and the twelve unit vectors outside $M$ are exactly those with $j$
negative. ($\varphi$ itself lies in $M$: it is
$-1 + 2\omega_1 + \omega_3 - 2\omega_1\omega_3$.) Where the distinction matters we write $M$ for
the index-$9$ subgroup; a reader coming from \cite{Engel} or \cite{Ruhland}
should read every statement here as being about the supergroup.

$\Lat$ is dense in the plane, and it is a \emph{group} and not a ring---of the
$900$ products of pairs of unit vectors only $684$ lie in it.
But $3\Lat$ is a ring, and is the maximal order of the field $\Lat$ spans:
\[
  \Lat = \tfrac13\,\Ok_V, \qquad V = \Q(\sqrt{33},\, i\sqrt3),
\]
which is Proposition~\ref{prop:order} and is what the rest of the paper runs on.
The \emph{Moser ring} $\Lat$ generates is strictly larger, with denominators
$3^k$; Gibbs' homomorphism and Hubai's $22$ colourings \cite{Hubai} live there,
and everything below is about $\Lat$ itself.

Parts \cite{Parts} classifies the record's constructions into four types, of
which \emph{type~M}---two copies of a base graph rotated relative to each other
and glued at a shared origin,
\[
  G = L \cup \theta(S), \qquad L, S \subseteq \Lat,
\]
---``is of the greatest interest, because so far it gives 5-chromatic graphs
with the least number of vertices''. The only rotation he records for it is
$\theta_4 = \arccos(7/8)$; writing in 2020 he could not have had the four that
\cite{Voronov} found the year after. Of the family he adds:
\begin{quote}
``Working constructions of type M with other rotations are not known yet.''
\end{quote}
That sentence is the question this paper takes up.

\emph{Two words, used throughout.} A rotation $\theta$ \textbf{joins} when there
are finite $L, S \subseteq \Lat$ for which $L \cup \theta(S)$, glued at the
shared origin, is not $4$-colourable---equivalently, by
Lemma~\ref{lem:subjoin}, when some such union contains an induced subjoin of
chromatic number exactly $5$. It is a property of the rotation and the lattice,
not of any pair of halves. A rotation is a \textbf{spindle} when its cosine is
rational; writing that cosine as $(2a-1)/(2a)$ names it $\theta_a$, and $\sqrt a$
is its \emph{shell} radius, the one radius at which a point and its image are a
unit apart. Section~\ref{sec:ball} fixes the normalisation that makes $a$ well
defined; $\theta_4$ is a spindle and is Heule's rotation.

\begin{theorem}[the divisibility law]\label{thm:law}
Let $a > \tfrac14$ be rational, so that $\theta_a$ is a rotation, and read
$4 \mid a$ as $v_2(a) \ge 2$---its meaning for integer $a$. Suppose the joint
lattice $\Lam_a = \Lat + \theta_a(\Lat)$ has rank~$8$---which by
Proposition~\ref{prop:rank} is every $a$ except those with $4a-1 = 3k^2$ or
$11k^2$ for a rational $k$. If $4 \nmid a$ then $\Lam_a$ admits a homomorphism
$4$-colouring. If $4 \mid a$ and the shell of $\Lat$ at radius $\sqrt a$ is
non-empty, it admits none.
\end{theorem}

A homomorphism $4$-colouring properly colours the \emph{entire} infinite lattice,
at every radius, so:

\begin{corollary}\label{cor:kill}
If $\Lam_a$ has rank $8$ and $4 \nmid a$, no subset of $\Lam_a$ is 5-chromatic;
in particular no 5-chromatic $L \cup \theta_a(S)$ exists, at any size.
\end{corollary}

The shell hypothesis carries only the second half, and it is sufficient rather
than necessary: what the proof needs is a vector of length $\tfrac12$ in
$\Lam_a$, which an occupied shell always supplies and an empty one may supply by
another route---$\Lam_{24}$ has twelve, exhibited in Section~\ref{sec:closed}.

\subsection*{What is proved here}

The engine is arithmetic, and it is the part likeliest to outlast the
obstruction. Proposition~\ref{prop:order} identifies $3\Lat$ as $\Ok_V$; from it,
Theorem~\ref{thm:primes} shows that $\Lat$'s two integral geometric
$4$-colourings are reduction modulo the two primes above $2$---a count due to
Dúcz, here given a structural reason---and Theorem~\ref{prop:shells} answers in
closed form which rationals are squared distances in $\Lat$: $a = n/b$ in lowest
terms is one exactly when $b \in \{1,3,9\}$ and $v_p(n)$ is even at every prime
splitting in $\Q(\sqrt{33})$ and inert in $\Q(\sqrt{-3})$. Corollary
\ref{cor:closed} makes the surviving integer rotations a set of density zero.
Proposition~\ref{prop:cross} removes the ball from the whole question: a cross
edge is a unit vector of $\Lam_a$, of which there are finitely many, exactly
enumerable.

Theorem~\ref{thm:law} is the application. It is one-sided and its reach is
bounded in a way worth stating at once: it kills every $a$ with $v_2(a) < 2$,
three integers in four, and
its own second half proves that at $4 \mid a$ with an occupied shell no
homomorphism colouring exists at all---so the method is silent exactly where the
working rotations live. It is silent outside the spindle family too, four of the
five rotations known to work having irrational cosine. Section
\ref{sec:stronger} kills eleven further integer candidates by a periodic
refinement, and Section~\ref{sec:seven} exhibits seven joining spindles, each
with a DRAT certificate a referee can check; Lemma~\ref{lem:chain} chains a
gadget of Parts' into a join at $a = 64n^2/9$ for every $n$, so infinitely many
spindles join.

\subsection*{What is not new}

The two geometric $4$-colourings of $\Lat$ and the fact that there are exactly
two are Dúcz's, after Hubai; his proof is a sublattice argument on coordinates
mod $2$, and contains no prime ideal and no number field, and his lattice is the
index-$9$ sublattice generated by $1, \omega_1, \omega_3, \omega_1\omega_3$, so
his count and Theorem~\ref{thm:primes} are about different groups. Parts already
writes his vertices over $\sqrt{33}$, subject to a congruence that normalises
coordinates inside one fixed lattice; the divisibility here is a condition on the
\emph{rotation}. The five previously known type-M rotations are Heule's and the
four of \cite{Voronov}. What is claimed below is the arithmetic, the law, and the
constructions of Section~\ref{sec:seven}.

That claim deserves a comparison, not a keyword search: the arithmetic framing
is not new even where the results are. The Polymath16
material \cite{Polymath} sets the Moser \emph{ring} inside
$\Q(\sqrt{-3},\sqrt{-11})$ and writes
it as $\Z[\tfrac{1+i\sqrt3}{2}, \tfrac{1+\sqrt{33}}{2}, \tfrac13]$; Speyer
\cite{Speyer} gives it as the elements integral over $\Z[\tfrac13]$, names the
ring $\Ok = \Z[\tfrac{1+\sqrt{33}}{2}]$ that Theorem~\ref{thm:primes} works over,
takes the norm $N(u) = u\bar u$ from one to the other, splits the ring modulo
$2^k$ into two factors and colours by one of the two projections. That splitting
at $k = 1$ is $\F_4 \times \F_4$, and $3$ is a unit modulo $2$, so both it and
the projections carry over to $\Lat$ unchanged: they are his, not ours. Gibbs
factors unit-modulus elements through $\Q[\sqrt{-p},\sqrt{-q}]$, and his
Theorem~1 \cite{GibbsRing} is the congruence description this paper works in:
the chart \eqref{eq:chart} with \eqref{eq:cong} is its $k = 0$ layer, scaled by
$3$. We claim none of that.

What we add is ideal-theoretic, and is what Proposition~\ref{prop:order} states:
that $3\Lat$ is the \emph{maximal} order of $V$, not merely an order inside it,
which is what makes the two factors the reductions modulo the two primes above
$2$ rather than a coordinate choice; that Theorem~\ref{thm:primes} counts them,
so that there are \emph{exactly} two and no third; and that occupancy of a shell
is a relative-norm condition with the closed-form answer of
Theorem~\ref{prop:shells}. Of the eight papers working on this
lattice---\cite{Ducz}, \cite{Parts}, \cite{Voronov}, \cite{Haugland},
\cite{Engel}, \cite{Ruhland}, \cite{Matolcsi},
\cite{DuczVarga}---none contains the words \emph{maximal order},
\emph{ring of integers}, \emph{prime ideal} or \emph{number field} anywhere in
its full text. The classical spectrum here is the triangular lattice's, the
Loeschian numbers $a^2+ab+b^2$; we have found no statement, in those papers or
in the Polymath16 material, of the corresponding set for the rank-$4$ lattice
$\Lat$, whose spectrum is Theorem~\ref{prop:shells}. That is evidence and not proof, and it is
evidence as searched in September 2026: this literature is partly distributed
across blog comments and data files---Speyer's is one---and a negative over it
cannot be certified.

\begin{figure}[ht]
\centering
\includegraphics{spindle.pdf}
\caption{The spindle at squared radius $a$.  Two shell points at distance $\sqrt a$
are carried onto each other by $\theta_a$ and land at distance exactly $1$, which
is what the rotation is for.  Half that displacement has length $\tfrac12$
always, since $\lvert 1-\theta_a\rvert = 1/\sqrt a$; whether it lies in the
lattice is the whole content of Theorem~\ref{thm:law}.  Drawn for $a=4$,
$\cos\theta_4 = 7/8$.}
\label{fig:spindle}
\end{figure}
\begin{figure}[ht]
\centering
\lawfig
\caption{Theorem~\ref{thm:law} as a sieve, over the integer spindles
$a \le 100$.  Top: the $59$ \emph{positive} squared radii at which $\Lat$ has points,
decided by the form \eqref{eq:norm} rather than by a ball, and in closed form by
Theorem~\ref{prop:shells}.  Middle: the multiples
of $4$.  Bottom: their intersection, the sixteen multiples of $4$ with an
\emph{occupied} shell---the candidates, in the sense Section~\ref{sec:seven}
fixes.  The nine gaps there are the multiples of $4$ whose shell is empty,
$a = 8$, $24$, $32$, $40$, $56$, $68$, $72$, $88$, $96$; they lie outside this
accounting altogether, because at an empty shell Theorem~\ref{thm:law} is silent
in both directions.  The law is a kill, not a construction, so the bottom
row also carries what became of each candidate: filled, the four given a verified
5-chromatic union in Section~\ref{sec:seven}; grey, the eleven the stronger
filter of Section~\ref{sec:stronger} kills; open, $a = 64$, on which that filter
degenerates at every modulus tried---it joins all the same, by
Lemma~\ref{lem:chain}, which builds a graph there rather than testing one.}
\label{fig:law}
\end{figure}
\section{The joint lattice, and why no ball is needed}\label{sec:ball}

\begin{proposition}\label{prop:cross}
Suppose $\Lam_a = \Lat \oplus \theta_a(\Lat)$ has rank $8$. Then
$(p,\theta_a(q)) \mapsto p - \theta_a(q)$ is a bijection from the cross edges of
$\Lat \cup \theta_a(\Lat)$---pairs $p \in \Lat$ and $\theta_a(q) \in
\theta_a(\Lat)$ at distance $1$ from each other---onto the unit vectors of
$\Lam_a$. There are finitely many, though the union is dense. Moreover such an edge is monochromatic under a
homomorphism $\varphi$ of $\Lam_a$ exactly when $\varphi(u) = 0$, so what a cross
edge contributes depends only on $u$.
\end{proposition}

\begin{proof}
If $\lvert p - \theta_a(q)\rvert = 1$ then $u = p - \theta_a(q)$ is a unit vector
of $\Lam_a$. Conversely the direct sum writes a unit vector uniquely as
$u = r + \theta_a(s)$, and then $p = r$, $q = -s$ is the only pair with
$p - \theta_a(q) = u$. Writing $\varphi = (\varphi_1,\varphi_2)$ on the two
summands, $\varphi(u) = \varphi_1(p) - \varphi_2(\theta_a(q))$, which vanishes
exactly when the two ends receive the same colour.
\end{proof}


\begin{proposition}\label{prop:rank}
Let $V \subset \C$ be the $\Q$-span of $\Lat$. Then $V = \Q(\sqrt{33},\,i\sqrt3)$,
a field of degree $4$ over $\Q$, and
\[
  \theta_a \in V
  \iff
  4a-1 = 3k^2 \text{ or } 11k^2 \text{ for a \emph{rational} } k.
\]
If $\theta_a \in V$ then $\Lam_a = \Lat + \theta_a(\Lat)$ has rank $4$; otherwise
it has rank exactly $8$ and Theorem~\ref{thm:law} applies. No intermediate rank
occurs.
\end{proposition}
\begin{proof}
The chart \eqref{eq:chart} with rational $A,B,C,D$ gives $V$ the $\Q$-basis
$1\mc \sqrt{33}\mc i\sqrt3\mc i\sqrt{11}$. Since $i\sqrt{11} = (i\sqrt3)\sqrt{33}/3$
and $(i\sqrt3)^2 = -3$, this is $\Q(\sqrt{33})(i\sqrt3)$, of degree $2$ over
$\Q(\sqrt{33})$ and $4$ over $\Q$; in particular $V$ is a \emph{field}.

Now $\theta_a = c + is$ with $c = (2a-1)/(2a) \in \Q$ and
$s = \sqrt{4a-1}/(2a)$, so $\theta_a \in V$ if and only if $is \in V$, that is
$\sqrt{4a-1} \in \Q\sqrt3 + \Q\sqrt{11}$. Writing
$\sqrt{4a-1} = p\sqrt3 + q\sqrt{11}$ and squaring gives
$4a-1 = 3p^2 + 11q^2 + 2pq\sqrt{33}$, so $pq = 0$ and $4a-1$ is $3p^2$ or
$11q^2$, with $p, q$ rational. When $a$ is an integer the rational $k$ is forced
to be one: if $4a-1 = 3(u/v)^2$ in lowest terms then $v^2 \mid 3$, so $v = 1$;
and then $k$ is odd besides, since $4a-1 \equiv 3 \pmod 4$ while $3k^2$ and
$11k^2$ are $0$ mod $4$ for even $k$. For non-integer $a$ the rational form is
the one that bites: $a = 5/9, 7/3, 13/3, 71/9, 31/3$ and $43/3$ all have
$4a-1 \in 3\Q^2 \cup 11\Q^2$, all have rank $4$, and all have a non-empty shell,
so they are genuinely excluded from Theorem~\ref{thm:law} and not vacuously. Both
branches occur: $7/3$, $13/3$, $31/3$, $43/3$ are $3k^2$ with $k \in \tfrac13\Z$,
while $5/9$ and $71/9$ are $11k^2$---$4a-1 = 11/9$ and $275/9$ at $k = 1/3$ and
$5/3$. An earlier version of this list ran only the $3k^2$ branch at non-integer
$a$ and so missed the second pair; it also included $1/3$, which is below the
normalisation of Section~\ref{sec:intro} and whose representative
$a/(4a-1) = 1$ is the integer already listed above.

If $\theta_a \in V$ then multiplication by it is a bijection of the field $V$, so
$\theta_a V = V$ and $\Lam_a \subset V$ has rank $4$. If $\theta_a \notin V$,
suppose $w \ne 0$ lay in $V \cap \theta_a V$, say $w = \theta_a v$ with
$v \in V \setminus \{0\}$. Then $\theta_a = wv^{-1} \in V$, because $V$ is a
field---a contradiction. So $V \cap \theta_a V = 0$ and
$\dim_\Q(V + \theta_a V) = 4 + 4 = 8$.
\end{proof}


\subsection*{One symmetry, and the normalisation it forces}
$\Lat$ is carried to itself by complex conjugation and by $-1$, so for any
$L, S \subseteq \Lat$,
\[
  \overline{L \cup \theta_a(S)} \;=\; \overline L \cup
  \theta_{a'}\bigl(-\overline S\bigr), \qquad a' = \frac{a}{4a-1},
\]
with $\overline L$ and $-\overline S$ again inside $\Lat$. The identity behind it
is $\theta_{a'} = -\overline{\theta_a}$, which one reads off
$\cos\theta_{a'} = -\cos\theta_a$ and $\sin\theta_{a'} = \sin\theta_a$. Mirror
images are isometric, so $\theta_a$ joins if and only if $\theta_{a'}$ does. The
map $a \mapsto a/(4a-1)$ is an involution exchanging $(\tfrac14,\tfrac12)$ with
$(\tfrac12,\infty)$ and fixing $\tfrac12$, and it leaves $v_2(a)$ alone whenever
the denominator of $a$ is odd.

\emph{We normalise to $\cos\theta_a \ge 0$, that is to $a \ge \tfrac12$}, and every
statement below about which $\theta_a$ join is a statement about such
representatives. This is not cosmetic. It is the shell that breaks the symmetry:
$\theta_4$ spins a shell of $30$ points, while its partner $\theta_{4/15}$---which
carries the mirror image of the very same $509$-vertex graph---has an empty shell,
since $15 \notin \{1,3,9\}$. Un-normalised, a criterion phrased through shells
would admit one of the pair and reject the other.

Then $\lvert p - \theta_a(p)\rvert = \lvert p\rvert \lvert 1-\theta_a\rvert =
\lvert p \rvert/\sqrt a$, so a point and its own image are a unit pair exactly
when $\lvert p\rvert = \sqrt a$: the rotation's own pairing lives on that
\emph{shell} and nowhere else, and Figure~\ref{fig:spindle} is the whole
mechanism in one picture. Other points do acquire cross edges---to images of
\emph{other} points, as Figure~\ref{fig:orbits} shows---but if the shell is empty
the rotation has nothing of its own to act on.

Voronov, Neopryatnaya and Dergachev \cite{Voronov} enumerate rotations carrying
one subset of vertices onto another; their first series runs on $\Lat$ and
reports five that work, Heule's among them, so five working type-M rotations were
known and not one. They are one orbit: all five are $\rho u$ with $\rho = \theta_4$
and $u$ a unit vector of $\Lat$, and exactly one of the five---$\rho$ itself---has
rational cosine. The other four are therefore not spindles, and in the family
classified here exactly one working member was known.


Every earlier test of a rotation built a ball of lattice points, found its unit
edges and ran a SAT solver. That has a free parameter---the ball---and a negative
answer about a ball is not an answer about the lattice: a rotation whose shell at
radius $\sqrt a$ lies outside the ball is being dismissed on a graph that never
contained a single pair $p, \theta_a(p)$. Section~\ref{sec:seven} exhibits six
rotations that a ball at the record's own scale misses for exactly that reason.

No ball is needed. Write $W$ for the $30$ unit vectors of $\Lat$. For a candidate
$\theta$ with $K = V(\theta)$ a CM field---every rotation considered here, since
$K = V(i\sqrt{4a-1})$ at a spindle and $V(\sqrt5)$ at the thirty $\rho u$ of
Section~\ref{sec:notclaimed}---both halves of $L \cup \theta(S)$ lie in
\[
  \Lam_\theta = \langle W \cup \theta(W) \rangle,
\]
the group they and their images generate: a finitely generated subgroup of
$(\C,+)$ with no radius in it. Its unit vectors are enumerable exactly:
$q(z) = \Tr_{K/\Q}(z\bar z)$ is a positive-definite quadratic form on
$\Lam_\theta \otimes \Q$, conjugation being the CM involution, and a unit vector
has $z\bar z = 1$ and hence $q(z) = \Tr(1)$, so the unit vectors lie \emph{among} the integer points of one
ellipsoid. That set is finite and can be listed exactly; sifting it by the exact
condition $z\bar z = 1$---which is four rational conditions where $q$ is one, or
one condition over $F$ where $q$ is a single rational---leaves the
unit vectors, a complete list obtained without reference to any ball.

\medskip
\noindent\textbf{The criterion.} Look for a group homomorphism
$\varphi : \Lam_\theta \to \mathcal{A}$ with $\lvert \mathcal{A} \rvert = 4$ and $\varphi(u) \ne 0$
for every unit vector $u$. If one exists then for any $p,q \in \Lam_\theta$ at
distance $1$ the difference $p-q$ is a unit vector, so $\varphi(p) \ne \varphi(q)$
and $p \mapsto \varphi(p)$ is a proper colouring of the whole infinite lattice.
Hence
\[
  \varphi \text{ exists} \implies \text{no subset of } \Lam_\theta \text{ is 5-chromatic, ever.}
\]
The converse fails---a lattice may be $4$-colourable by a colouring that is not a
homomorphism---so the test is a \emph{one-sided kill}, and a survivor is a
candidate and nothing more. There are two groups of order $4$, the Klein group and
$\Z/4$, and both must be ruled out.

A one-sided kill has one way to be wrong: killing something that works. Every
rotation known to carry a 5-chromatic type-M graph is therefore a test of the
criterion, and there are five---Heule's and the four of \cite{Voronov}. Run
against all five, the criterion returns \emph{unknown} each time and kills none;
the five joint lattices have rank $8$ and carry $126$, $108$, $108$, $90$ and
$90$ unit vectors respectively. That is the only outcome consistent with their
being 5-chromatic.

Four of the five are informative and one is not. Heule's rotation is
$\theta_4$, and $4 = 4\cdot1^2$, so Proposition~\ref{prop:free} forces its
survival on any lattice closed under $\omega$: $\Lam_4$ holds a half-unit vector
and the criterion could not have killed it whatever the implementation did. The
test is therefore a test on the other four, where nothing forces the outcome.


\section{The criterion in closed form}\label{sec:closed}

\begin{proposition}\label{prop:free}
Let $\Lam$ be any subgroup of $(\C,+)$ closed under $\omega = e^{i\pi/3}$ and
containing a unit vector $g$, and let $a = 4m^2$ for a positive integer $m$. Then
$\Lam + \theta_a(\Lam)$ contains a vector of length $\tfrac12$, so by
Proposition~\ref{prop:half} the criterion cannot kill $\theta_a$: neither group
of order $4$ admits a homomorphism avoiding the unit vectors. Both lattices here
are $\omega$-closed---$\Lat$ by construction, and Haugland's because $\omega$ is
one of its $42$ roots of unity.
\end{proposition}

\begin{proof}
$\lvert 2mg\rvert = 2m = \sqrt a$, so $2mg$ lies on the shell, and
\[
  h_g \;=\; \tfrac12\bigl(2mg - \theta_a(2mg)\bigr) \;=\; m(1-\theta_a)g
\]
lies in $\Lam + \theta_a(\Lam)$ with length $m \lvert 1-\theta_a\rvert =
m/\sqrt a = \tfrac12$. Proposition~\ref{prop:half} does the rest, both halves of
it, since $\Lam$ is $\omega$-closed.
\end{proof}

\begin{proposition}\label{prop:half}
Suppose $\Lam_\theta$ contains a vector of length $1/2$. Then no homomorphism
onto a group of order $4$---Klein or cyclic---avoids every unit vector.
\end{proposition}

\begin{proof}
Let $\lvert h \rvert = 1/2$, so $2h$ is a unit vector of $\Lam_\theta$. In the
Klein group every element satisfies $2x = 0$, so $\varphi(2h) = 2\varphi(h) = 0$
and the unit vector $2h$ is killed.

In $\Z/4$, $\varphi(2h) = 2\varphi(h) \ne 0$ forces $\varphi(h)$ odd, and that is
already impossible. Write $\omega = e^{i\pi/3}$, a unit vector of $\Lat$ with
$\omega\Lat = \Lat$, so $\omega\Lam_\theta = \Lam_\theta$ and $\omega h$,
$\omega^2 h$ are half-unit vectors alongside $h$. Since $\omega^2 = \omega - 1$,
\[
  h + \omega h + \omega^2 h \;=\; 2\omega h \;\in\; 2\Lam_\theta ,
\]
so $\varphi(h) + \varphi(\omega h) + \varphi(\omega^2 h) = 2\varphi(\omega h)$ is
even, while three odd values sum to an odd one.
\end{proof}

Both halves are unconditional, and the $\omega$-identity is why: it is a relation
of \emph{odd} weight among half-unit vectors, and it exists wherever one
half-unit vector does. The sweep runs over every multiple of $4$ up to
$a = 100$, and those are the only integer spindles that can have a half-unit
vector at all: where $4 \nmid a$, Theorem~\ref{thm:law} supplies a homomorphism
$4$-colouring, which Proposition~\ref{prop:half} forbids wherever one exists.
Nineteen of them do have one---the sixteen with an occupied shell, where
Theorem~\ref{thm:law}'s own construction produces it, together with $a = 24$,
$72$ and $88$, where the shell is empty and twelve appear regardless---and on
every one of the nineteen the relation is found, as it must be.


Half-unit vectors are therefore the object to count. The shell supplies a family
of them in closed form---though not, as $a = 24$ will show below, all of them.
For $p$ on the shell, $p$ and $\theta_a(p)$ are both at distance
$\sqrt a$, and
\[
  \tfrac12\big(p - \theta_a(p)\big)
\]
has length exactly $1/2$ for every $a$, since $\lvert 1 - \theta_a \rvert =
1/\sqrt a$. Now $p - \theta_a(p)$ always lies in $\Lam_a$, being a difference of
one element of $\Lat$ and one of $\theta_a(\Lat)$; whether \emph{half} of it does
is another matter, and that is the whole content of the classification:
$\tfrac12(p-\theta_a(p)) \in \Lam_a$ exactly when $p - \theta_a(p) \in 2\Lam_a$.

That classifies the half-unit vectors \emph{the shell produces}, and it is not a
classification of all of them. At $a = 24$ the shell is empty, so the
construction above yields nothing at all---and yet $\Lam_{24}$ holds twelve
vectors of length $\tfrac12$. One of them needs no machinery to write down. Since
$\Lat \cap \Q = \tfrac13\Z$, both $r = -\tfrac53$ and $s = 2$ lie in $\Lat$, and
with $\cos\theta_{24} = \tfrac{47}{48}$, $\sin\theta_{24} = \tfrac{\sqrt{95}}{48}$,
\[
  r + \theta_{24}(s)
  \;=\; -\tfrac53 + 2\Big(\tfrac{47}{48} + i\,\tfrac{\sqrt{95}}{48}\Big)
  \;=\; \frac{7 + i\sqrt{95}}{24},
  \qquad
  \Big\lvert \frac{7+i\sqrt{95}}{24} \Big\rvert^2 = \frac{49+95}{576} = \frac14 .
\]
So Proposition~\ref{prop:half} applies at $a = 24$ and no homomorphism
$4$-colouring exists there, although Theorem~\ref{thm:law} says nothing: its
second half asks for an occupied shell, and an occupied shell is sufficient for a
half-unit vector without being necessary. In the language of Step~1 below this
is the case $s = \lambda r$ with $\lambda = -\tfrac65$, where
$\lvert r \rvert^2\big(a(1+\lambda)^2 - \lambda\big)/a = \tfrac14$; the shell is
$\lambda = -1$, and it is not the only value that can reach length $\tfrac12$.
Section~\ref{sec:stronger} takes this up again, where it explains survivals the
law cannot.

\section{The arithmetic of the Moser lattice}\label{sec:arith}

\subsection{A module over a real quadratic ring}
In the chart
\begin{equation}\label{eq:chart}
  z = \frac{A + B\sqrt{33}}{12} + i\,\frac{C\sqrt3 + D\sqrt{11}}{12},
  \qquad A,B,C,D \in \Z,
\end{equation}
the lattice $\Lat$ is cut out by the congruences
\begin{equation}\label{eq:cong}
  \text{(i)}\quad A \equiv B \equiv C \equiv D \pmod 2,
  \qquad
  \text{(ii)}\quad A + B + C + D \equiv 2A \pmod 4,
\end{equation}
sixteen of the $256$ residue classes mod $4$. (The chart without them is a group
containing $\Lat$ with index $16$, so reading a chart count as a lattice count
overcounts---at rational squared radii at least; Lemma~\ref{lem:auto} below says
it never does at integer ones.) Squared lengths are
\begin{equation}\label{eq:norm}
  \lvert z \rvert^2
  = \frac{\big(A^2 + 33B^2 + 3C^2 + 11D^2\big) + 2(AB + CD)\sqrt{33}}{144},
\end{equation}
so they lie in $\Q(\sqrt{33})$, and are \emph{rational exactly when
$AB + CD = 0$}, in which case they lie in $\tfrac1{144}\Z$. Write
\[
  S = A^2 + 33B^2 + 3C^2 + 11D^2
\]
for the numerator, so a rational squared length is $S/144$. (This $S$ is the
form's value throughout Sections \ref{sec:arith} and \ref{sec:proof}; the small
half of a construction carries the same letter only outside them.)

At \emph{integer} squared radii the congruences cost nothing, which is worth
recording because it is what makes the chart safe to compute with there.

\begin{lemma}\label{lem:auto}
Let $A,B,C,D \in \Z$ satisfy $AB + CD = 0$ and $16 \mid S$. Then \eqref{eq:cong}
holds automatically.
\end{lemma}

\begin{proof}
Squares modulo $16$ depend only on their bases modulo $8$, so both hypotheses and
both conclusions are determined by $(A,B,C,D)$ modulo $8$, and $8^4 = 4096$
classes decide the statement. Of those, $320$ satisfy $16 \mid S$ while failing
\eqref{eq:cong}, and not one of the $320$ has $AB+CD \equiv 0 \pmod 8$
(\texttt{audit\_numbers.py}).
\end{proof}

If $\lvert z\rvert^2 = a$ is a positive integer then $S = 144a$, so $16 \mid S$
and $AB + CD = 0$, and Lemma~\ref{lem:auto} applies: every integral chart point
at an integer squared radius already lies in $\Lat$. Measured, the two agree
exactly---identical counts at every $a \le 300$, not merely the same emptiness.

The module structure is in fact a great deal more than a module structure.

\begin{proposition}\label{prop:order}
$3\Lat$ is the maximal order $\Ok_V$ of $V = \Q(\sqrt{33}, i\sqrt3)$; that is,
$\Lat = \tfrac13\Ok_V$. Writing $F = \Q(\sqrt{33})$ for the real quadratic
subfield and $\bar{\phantom{z}}$ for complex conjugation, which generates
$\Gal(V/F)$,
\[
  \lvert z \rvert^2 = z\bar z = \Nm_{V/F}(z),
\]
so a squared radius of $\Lat$ is one ninth of a relative norm of an algebraic
integer.
\end{proposition}

\begin{proof}
Write coordinates over the $\Q$-basis $(1, \sqrt{33}, i\sqrt3, i\sqrt{11})$ of
$V$. Reducing the $30$ generators of $\Lat$ to Hermite form and multiplying by
$3$ gives a $\Z$-basis of $3\Lat$:
\[
  b_1 = \tfrac14(1,1,1,-1), \dsep
  b_2 = \tfrac12(0,1,0,-7), \dsep
  b_3 = \tfrac12(0,0,1,1), \dsep
  b_4 = (0,0,0,1).
\]
It contains $1$, because $\tfrac13 \in \Lat$ by \eqref{eq:chart} at $(4,0,0,0)$,
and it is closed under multiplication: the sixteen products $b_ib_j$ reduce to
integer combinations of the $b_k$ by the rules $\sqrt3\sqrt{11} = \sqrt{33}$,
$\sqrt3\sqrt{33} = 3\sqrt{11}$ and $\sqrt{11}\sqrt{33} = 11\sqrt3$. So $3\Lat$ is
an order in $V$ (\texttt{moser\_is\_ov.py}). (The same closure holds on all $900$
products of unit vectors, where $\Lat$ itself fails: only $684$ of them return.)

For the discriminant, $\Tr_{V/\Q}(x) = 4x_1$ on these coordinates, since the four
embeddings negate $\sqrt{33}$ and $i\sqrt3$ independently and kill every other
basis vector. Hence
\[
  \bigl(\Tr(b_ib_j)\bigr) =
  \begin{pmatrix}
      5 & -22 &   4 &  11\\
    -22 &-506 &  77 & 154\\
      4 &  77 & -14 & -22\\
     11 & 154 & -22 & -44
  \end{pmatrix},
  \qquad \det = 1089 .
\]
Now compare with the field. Put $F = \Q(\sqrt{33})$, so $V = F(i\sqrt3)$ and
$\operatorname{disc}(F) = 33$ because $33 \equiv 1 \pmod 4$. Two standard facts
close the argument without computing $\operatorname{disc}(V)$ separately. First,
for any order $\mathcal O$ of $V$,
\[
  \operatorname{disc}(\mathcal O) = [\Ok_V : \mathcal O]^2 \operatorname{disc}(V).
\]
Second, the tower formula for relative discriminants, applied to $V/F$ of
degree $2$ gives
$\operatorname{disc}(V) = \Nm_{F/\Q}(\mathfrak d_{V/F})\operatorname{disc}(F)^2$,
so $\operatorname{disc}(V)$ is a positive multiple of $33^2 = 1089$. Taking
$\mathcal O = 3\Lat$ and substituting,
\[
  1089 \;=\; [\Ok_V : 3\Lat]^2 \cdot \Nm_{F/\Q}(\mathfrak d_{V/F}) \cdot 1089 ,
\]
which forces both factors to be $1$. So $3\Lat = \Ok_V$, and as a by-product
$\mathfrak d_{V/F} = (1)$: $V/F$ is unramified at every finite prime, which is
what Theorem~\ref{prop:shells} will use. Remark~\ref{rem:narrow} reaches the same
conclusion from genus theory, and identifies $V$ as the narrow Hilbert class
field of $\Q(\sqrt{33})$.

The norm identity is the definition of $\Nm_{V/F}$: complex conjugation fixes
$\sqrt{33}$ and negates $i\sqrt3$, so it generates $\Gal(V/F)$ and
$z\bar z = \Nm_{V/F}(z)$.
\end{proof}

Two consequences follow, and one of them is the whole of Section
\ref{sec:seven}'s criterion.
\begin{itemize}
  \item $\Lat$ \emph{is a module over} $\Ok = \Z[(1+\sqrt{33})/2]$, the ring of
  integers of $\Q(\sqrt{33})$. This is a finite check and not a sample. The
  congruences \eqref{eq:cong} depend only on $(A,B,C,D)$ modulo $4$, and exactly
  $16$ of the $256$ classes satisfy them. Multiplication by $\sqrt{33}$ sends
  $(A,B,C,D)$ to $(33B,\,A,\,11D,\,3C)$---using $\sqrt3\sqrt{33} = 3\sqrt{11}$
  and $\sqrt{11}\sqrt{33} = 11\sqrt3$---and carries each of the $16$ classes to
  one of the $16$. For $(1+\sqrt{33})/2$ the same map, halved, is integral
  because $A \equiv B$ and $C \equiv D \pmod 2$ on those classes, and again
  preserves them. Since $\rank_\Z \Lat = 4$ and $\rank_\Z \Ok = 2$, the lattice
  $\Lat$ has rank $2$ over $\Ok$.
  \item $\Lat$ \emph{is not a lattice} in the geometric sense. It is dense, for
  the reason given in Section~\ref{sec:intro}; its short vectors are as short as one likes, and a
  depth-$4$ ball already contains one with
  $\lvert v \rvert^2 = (23 - 4\sqrt{33})/9 \approx 0.0024$. Any argument wanting a
  smallest vector must restrict to \emph{rational} norms, where
  $\tfrac1{144}\Z$ supplies one.
\end{itemize}

The smallest rational squared norms of $\Lat$, decided exactly from \eqref{eq:norm}
together with \eqref{eq:cong}, are
\[
  \tfrac19,\ \tfrac13,\ \tfrac49,\ \tfrac59,\ \tfrac79,\ 1,\ \dots
\]
so $1/9$ is the least, attained by six points, and it is the only one below
$1/4$. The value $1/36$ is not attained at all, a fact used in Section
\ref{sec:proof}.

\subsection{The two colourings are prime reductions}
Since $33 \equiv 1 \pmod 8$, the prime $2$ splits in $\Q(\sqrt{33})$:
\[
  2 = -\pi_1\pi_2,
  \qquad
  \pi_1 = \frac{5+\sqrt{33}}{2},
  \qquad
  \pi_2 = \frac{5-\sqrt{33}}{2},
  \qquad
  \Nm(\pi_i) = -2 .
\]
Because $\Lat$ is an $\Ok$-module, $\pi_i\Lat$ is a subgroup, and
$[\Lat : \pi_i \Lat] = \Nm(\pi_i)^2 = 4$, exactly the index a $4$-colouring needs,
where $[\Lat : 2\Lat] = 16$ is four times too big.

\begin{theorem}[the colourings are prime reductions]\label{thm:primes}
The rational prime $2$ splits in $\Q(\sqrt{33})$ as $-\pi_1\pi_2$ with
$\pi_{1,2} = (5 \pm \sqrt{33})/2$. The lattice $\Lat$ is a module over
$\Ok = \Z[(1+\sqrt{33})/2]$. Call a proper $4$-colouring of $\Lat$
\emph{integral} when its classes are cosets of a subgroup and \emph{geometric}
when $\Lat$'s rotations permute them. The two such colourings are the reductions
modulo $\pi_1\Lat$ and $\pi_2\Lat$: two kernels, so two colourings up to
permutation of the colours, and two because $2$ splits.
\end{theorem}

\begin{proof}[Proof of Theorem \ref{thm:primes}]
First, the quotient has no choice but to be Klein. By integrality the colour
classes are the cosets of a subgroup $H$ of index $4$. By geometry the isometry
$\omega = e^{i\pi/3}$ of $\Lat$ permutes those classes, and it fixes the origin,
so it fixes the class containing it: $\omega H = H$, and $\omega$ induces an
automorphism $T$ of $\Lat/H$. Since $\omega^2 = \omega - 1$,
\[
  T^2 - T + I = 0 .
\]
If $\Lat/H$ were cyclic of order $4$ then $T = \pm I$, the only automorphisms it
has, and the left side would be $I$ or $3I$---neither zero in $\operatorname{End}(\Z/4)$. So
$\Lat/H$ is the Klein group, $2\Lat \subseteq H$, and the search may be confined to
subgroups containing $2\Lat$.

That done, a homomorphism onto the Klein group killing no unit vector is the same
thing as a subgroup of index $4$ containing $2\Lat$ and meeting none of the $30$
unit vectors, so the whole question lives in $\Lat/2\Lat$---and
Proposition~\ref{prop:order} says exactly what that quotient is. Multiplication by
$3$ carries $\Lat$ isomorphically onto $\Ok_V$ and $\pi_i\Lat$ onto
$\pi_i\Ok_V$, so $\Lat/\pi_i\Lat = \Ok_V/\pi_i\Ok_V$; and $\pi_i$ is inert in
$V$, since the completion of $F$ at $\pi_i$ is $\Q_2$---the prime $2$ splits---and
$-3 \equiv 5 \pmod 8$ generates its unramified quadratic extension. Each quotient
is therefore the field $\F_4$, and as $\pi_1, \pi_2$ are coprime the Chinese
remainder theorem gives a ring isomorphism of $\Ok_V/2\Ok_V$ onto
$\F_4 \times \F_4$, which multiplication by $3$ transports to an isomorphism of
\emph{additive groups}---$\Lat$ is not a ring, and only the additive structure is
used below:
\[
  \Lat/2\Lat \;\xrightarrow{\ \sim\ }\; \Lat/\pi_1\Lat \times \Lat/\pi_2\Lat
  \;=\; \F_4 \times \F_4 ,
\]
under which the two prime reductions are the coordinate lines:
$\pi_1\Lat/2\Lat = 0 \times \F_4$ and $\pi_2\Lat/2\Lat = \F_4 \times 0$.

Now the unit vectors. A unit vector $u$ has $\lvert u \rvert^2 = 1$, so
$v_{\pi_i}(u) = 0$ at both primes by Lemma \ref{lem:val}: its class is invertible
in each factor. There are $3 \cdot 3 = 9$ such classes, and all nine occur. The
thirty unit vectors are the $\omega^k\varphi^j$ with $0 \le k < 6$ and
$\lvert j \rvert \le 2$, for the same $\omega = e^{i\pi/3}$ and the half-angle
$\varphi = e^{i\theta_3/2}$ of Section~\ref{sec:intro}; modulo $2\Lat$,
$\omega^3 = -1 \equiv 1$ and $\varphi^{-1} \equiv \varphi^2$, so their classes are
the nine with $k, j \in \{0,1,2\}$, and those nine are distinct. So the $9$ is a ceiling
attained, not a coincidence of a count.

Two lines finish the theorem. The image $W$ of an index-$4$ subgroup containing
$2\Lat$ is an $\F_2$-subspace of $\F_4 \times \F_4$ of dimension $2$. If it meets
no unit class then every element of $W$ has a zero coordinate, so
\[
  W \;=\; \big(W \cap (\F_4 \times 0)\big) \;\cup\; \big(W \cap (0 \times \F_4)\big) ,
\]
a union of two subgroups---and a group is never the union of two proper subgroups.
Hence $W$ lies inside one coordinate line, and both lines have $\F_2$-dimension
$2$, so $W$ \emph{is} that line. Exactly two subgroups survive, and they are the
prime reductions:
\[
  \ker\gamma_1 = \pi_1\Lat, \qquad \ker\gamma_2 = \pi_2\Lat .
\]
The count---two---is the one D\'ucz reaches for his own lattice, by enumeration
rather than by his sublattice argument, and reached here for $\Lat$, which is not
his lattice.

Both are geometric, which closes the circle: by Lemma \ref{lem:val},
$v_{\pi_i}(\lvert x\rvert^2) = 2v_{\pi_i}(x)$ for $x \ne 0$, so $x \in \pi_i\Lat$
if and only if $v_{\pi_i}(\lvert x \rvert^2) \ge 2$---membership in either kernel
is read off the squared length alone, and monochromatic pairs are decided by
distance.
\end{proof}

In coordinates, for checking: take the Hermite basis of $\Lat$ in the chart
\eqref{eq:chart}, rows in the order $(A,B,C,D)$,
\[
  b_1 = (1,1,1,-1), \quad b_2 = (0,2,0,-14), \quad
  b_3 = (0,0,2,2), \quad b_4 = (0,0,0,4),
\]
and write a class of $\Lat/2\Lat \cong \F_2^4$ as the integer whose bit $i$ is the
coefficient of $b_{i+1}$. Then
\[
  \pi_1\Lat/2\Lat = \{0,5,9,12\}, \qquad \pi_2\Lat/2\Lat = \{0,3,4,7\},
\]
the nine unit classes are the nine integers in neither list, and a sweep over all
$35$ two-dimensional subspaces of $\F_2^4$---the enumeration the argument above
replaces---returns those two and no others (\texttt{two\_colourings.py}). (The
masks depend on the basis and the bit order; both are fixed here so the numbers
can be checked.)

Stability is the wrong invariant, and it is worth saying which one fails, because
matching the $\Ok$-stable subgroups against D\'ucz's two is the natural first move
and cannot work. Multiplication by $\sqrt{33}$ is the \emph{identity} on
$\Lat/2\Lat$, since $33 \equiv 1$, so it constrains nothing; multiplication by
$(1+\sqrt{33})/2$ leaves eleven of the thirty-five two-dimensional subspaces
stable---the two eigenspaces and the nine sums of a line from each---not two. It
is avoiding the unit vectors that cuts the field to two, and the ring structure
above that names them.

The theorem counts the Klein colourings, and the empty-shell verdicts of
Section~\ref{sec:stronger} need the cyclic ones too, so count them here. A homomorphism $\Lat \to \Z/4$ missing every
unit vector is determined by its values on a basis; writing those as
$\ell + 2s$ over $\F_2$ makes the condition one affine system per $\ell$, a sweep
over $2^4$ rather than $4^4$, exactly as in Proposition~\ref{prop:half}'s proof.
The answer is small enough to print. In the Hermite basis $b_1,\dots,b_4$ fixed
below, writing a homomorphism as the tuple $(\varphi(b_1),\dots,\varphi(b_4))$,
exactly twelve of the $256$ tuples kill no unit vector:
\[
  (0,0,2,1),\quad (0,1,0,0),\quad (1,0,1,3),\quad
  (1,1,0,0),\quad (1,1,2,1),\quad (1,3,1,3)
\]
and their negatives. Up to negation---the only non-trivial automorphism of
$\Z/4$---there are therefore exactly \emph{six}. So $\Lat$ carries eight homomorphism $4$-colourings altogether, two
Klein and six cyclic, each counted up to the automorphisms of its own group, and
it is that $8$ that the interface argument pairs with itself; that computation
is \texttt{scripts/interface\_matchings.py} in the archive.

The colour condition is therefore divisibility in $\Ok$,
\begin{equation}\label{eq:colour}
  \gamma_i(x) = 0 \iff \pi_i \mid x ,
\end{equation}
proved above and confirmed directly, point by point and for both colourings, on a
ball of the lattice; the accompanying code reports the size. The number of
colourings is two \emph{because $2$ splits into exactly two primes}: the subgroups
that avoid every unit class are the two prime reductions, and there are two of
those because $33 \equiv 1 \pmod 8$. That is a reason, and D\'ucz's sublattice
proof does not give one. Note also
$\pi_1\Lat \cap \pi_2\Lat = \pi_1\pi_2\Lat = 2\Lat$: both colours vanish exactly
on $2\Lat$.

\section{Proof of the divisibility law}\label{sec:proof}

Assume $\Lam = \Lat + \theta_a(\Lat)$ has rank $8$, so
$\Lam = \Lat \oplus \theta_a(\Lat)$ and a homomorphism $\varphi : \Lam \to \mathcal{A}$ is a
pair $(\varphi_1, \varphi_2)$, one per summand, and each $\varphi_i$ must avoid
the $30$ unit vectors of its own summand.

Take $\mathcal{A}$ to be the Klein group first. Its exponent is $2$, so
$2\Lat \subseteq \ker\varphi_i$ automatically---no geometric hypothesis is needed
and none is available, since a homomorphism avoiding the unit vectors need not be
geometric: Section~\ref{sec:arith} exhibits six cyclic ones that are not. What is
available is the enumeration inside Theorem~\ref{thm:primes}: a subgroup of index
$4$ containing $2\Lat$ and meeting none of the nine unit classes is $\pi_1\Lat$ or
$\pi_2\Lat$. The image cannot be smaller than the whole Klein group either. Were
it of order $2$, then $1 + \omega + \omega^2 = 2\omega$ would give
$\varphi_i(1) + \varphi_i(\omega) + \varphi_i(\omega^2) = \varphi_i(2\omega) = 0$,
while each of the three terms is the one non-zero element of an order-$2$ group
and their sum is that element again. So each $\varphi_i$ is reduction modulo
$\pi_1\Lat$ or $\pi_2\Lat$, and with the six automorphisms of the Klein group that
is at most $2 \cdot 2 \cdot 6 = 24$ candidates, whatever $a$ is. (Both reductions
happen to be geometric, by Theorem~\ref{thm:primes}, but that is a conclusion
here and not a hypothesis.)

The cyclic target $\mathcal{A} = \Z/4$ is disposed of once and for all by
Proposition~\ref{prop:half}, which Step 5 invokes; the six cyclic colourings
counted in Section~\ref{sec:arith} are the ones that survive on $\Lat$ alone and
play no part here. The question is which of
them a \emph{cross} unit vector $u = r + \theta(s)$, with $r,s \in \Lat$ both
non-zero, destroys. Figure~\ref{fig:orbits} shows the two orbits these fall into.

The five steps deliver, in order: that a cross unit vector has its two parts
parallel, $s = \lambda r$ with $\lambda$ a real scalar (Step 1); that
$\lambda = -1$ is the shell, the only case carrying geometric information
(Step 2); that a shell point lies in $2\Lat$ exactly when $4 \mid a$, which is
where the divisibility enters (Step 3); that every other $\lambda$ constrains
nothing the shell has not already constrained (Step 4); and that the two
conclusions follow---a colouring exists when $4 \nmid a$, and Proposition
\ref{prop:half} forbids one when $4 \mid a$ with the shell occupied (Step 5).

\subsection*{Step 1: the two parts are parallel}
Writing $w = r\bar s$,
\[
  \lvert u \rvert^2
  = \lvert r \rvert^2 + \lvert s \rvert^2
  + \frac{(2a-1)\operatorname{Re} w + \sqrt{4a-1}\operatorname{Im} w}{a}.
\]
Separate the coefficients over $F = \Q(\sqrt{33})$. By \eqref{eq:norm} the
squared norms $\lvert r \rvert^2$ and $\lvert s \rvert^2$ lie in $F$, and so does
$\operatorname{Re} w$; the chart's imaginary part is $(C\sqrt3 + D\sqrt{11})/12$
and $\sqrt{11} = \sqrt{33}/\sqrt3$, so $\operatorname{Im} w$ lies in
$\sqrt3\,F$. Then $\lvert u \rvert^2 = 1$ reads
\[
  \alpha + \beta\sqrt{3(4a-1)} = 1, \qquad \alpha, \beta \in F,
\]
and rank $8$ is exactly the statement that $\sqrt{3(4a-1)} \notin F$---which by
Proposition~\ref{prop:rank} fails only at $4a-1 = 3k^2$ or $11k^2$. Note this is
the $\Q$-span and not the field: at $a = 1/2$ the root $\sqrt{4a-1} = 1$ lies in
$\Q(\sqrt3,\sqrt{11})$ and the rank is $8$ regardless. So $\beta = 0$, the
$\sqrt{4a-1}$ term vanishes on its
own: $\operatorname{Im} w = 0$, that is, $s = \lambda r$ for a real $\lambda$.

\begin{figure}[ht]
\centering
\includegraphics{orbits.pdf}
\caption{The entire interface of $\theta_4$: all $126$ cross edges a depth-$4$ ball
offers, computed from the coordinates.  Sixty run to the shared origin---the
inner star---thirty join two points of the shell at radius $2$, and thirty-six
run crosswise between the dotted radii $\tfrac{\sqrt{11}}{2}\mp\tfrac{\sqrt3}{6}
= 1.3696$ and $1.9470$.  Both halves here are the same ball. One edge of each non-trivial
orbit is drawn heavy, its end in $L$ filled and its end in $\theta_4(L)$ open.  Every edge in the figure
has length exactly $1$.}
\label{fig:orbits}
\end{figure}


\subsection*{Step 2: $\lambda = -1$ gives the shell}
Then $u = (1-\theta)r$, and $\lvert 1 - \theta \rvert = 1/\sqrt a$ forces
$\lvert r \rvert^2 = a$: the point $r$ lies on the shell and $u$ is a shell
displacement. Since $s = -r$ and the colour group has exponent $2$,
$\gamma(s) = \gamma(r)$: the pair of colours is \emph{diagonal}, $(c,c)$, and
$\varphi(u) = \psi_1(c) + \psi_2(c)$, where $\psi_1,\psi_2$ are the two
automorphisms---not the two primes, which keep the letter $\pi$ throughout. The
colour group has exponent $2$, so $\varphi(u) = 0$ exactly when
$\psi_1(c) = \psi_2(c)$. At $c = 0$ that holds for \emph{every} candidate, and
$c = 0$ is exactly $r \in 2\Lat$, both primes at once. For $c \ne 0$ it still
holds for those candidates whose two automorphisms happen to agree at $c$, which
is why the conclusion below asks for a relative $3$-cycle and not merely for
$c \ne 0$.

\subsection*{Step 3: $r \in 2\Lat \iff 4 \mid a$}
Here $S = 144a$ is an integer even when $a$ is not, because $A,B,C,D$ are.
Then $v_2(S) = 4 + v_2(a)$, so $4 \mid a$---that is
$v_2(a) \ge 2$---is exactly $64 \mid S$. Recall too that $\lvert r \rvert^2$ is
rational exactly when $AB + CD = 0$. Three lemmas settle it.

\begin{lemma}\label{lem:A}
For $(A,B,C,D)$ satisfying \eqref{eq:cong} with $AB + CD = 0$ we have
$16 \mid S$. Equivalently, every rational squared norm of $\Lat$ lies in
$\tfrac19\Z$.
\end{lemma}

\begin{proof}
For odd $x$ write $x^2 = 1 + 8k_x$; then $x \mapsto k_x \bmod 2$ is the
homomorphism $\varepsilon : (\Z/8)^\times \to \Z/2$ with $\varepsilon = 0$ on
$\pm1$ and $1$ on $\pm3$. If all four are odd then
$S \equiv 8\big(\varepsilon(A)+\varepsilon(B)+\varepsilon(C)+\varepsilon(D)\big)
\pmod{16}$, and $AB = -CD$ gives
\[
  \varepsilon(A)+\varepsilon(B) = \varepsilon(AB) = \varepsilon(-CD)
  = \varepsilon(C)+\varepsilon(D),
\]
so the bracket is even. If all four are even then $S = 4S'$ with
$S' = A'^2 + 33B'^2 + 3C'^2 + 11D'^2$, and mod $4$,
$S' \equiv [A'] + [B'] + 3[C'] + 3[D']$ with $[\,\cdot\,]$ the odd indicator; the
congruence (ii) makes the number of odd entries $0$, $2$ or $4$: halving
\eqref{eq:cong} gives $2(A'+B'+C'+D') \equiv 4A' \pmod 4$, so
$A'+B'+C'+D'$ is even. Of the six
two-element patterns $\{A',B'\}$ and $\{C',D'\}$ are excluded by
$A'B' + C'D' = 0$ (one product odd, the other even), while the rest give
$1 + 3 \equiv 0$.
\end{proof}

\begin{lemma}\label{lem:B}
If $A,B,C,D$ are all odd and $AB + CD = 0$ then $S \equiv 16 \pmod{32}$. The
second hypothesis is not decorative: $(1,1,1,3)$ gives $S = 136 \equiv 8$.
\end{lemma}

\begin{proof}
Refine $\varepsilon$ one step: for odd $x$ write $x^2 = 1 + 8k_x$ as before and
set $\eta(x) = k_x \bmod 4$, so $x^2 \equiv 1 + 8\eta(x) \pmod{32}$. From
$AB = -CD$ we get $A^2B^2 = C^2D^2$, whose two sides are $1 + 8(\eta_A+\eta_B)$
and $1 + 8(\eta_C+\eta_D)$ modulo $32$, the cross terms $64\eta\eta$ vanishing;
so $\eta_A + \eta_B \equiv \eta_C + \eta_D \pmod 4$. Expanding $S$ modulo $32$
gives $16 + 8(\eta_A + \eta_B - \eta_C - \eta_D) \equiv 16$.
\end{proof}

\begin{lemma}\label{lem:C}
If $4 \mid X^2 + 33Y^2 + 3Z^2 + 11W^2$ then $X+Y+Z+W$ is even.
\end{lemma}

\begin{proof}
Modulo $4$ a square is the odd indicator, so the form is congruent to
$[X] + [Y] + 3[Z] + 3[W]$. One odd entry gives $1$ or $3$; three give
$5 \equiv 1$ or $7 \equiv 3$. Neither is $0$, so the number of odd entries is
even.
\end{proof}

\emph{Forward.} If $r = 2w$ with $w \in \Lat$ then $a = 4\lvert w \rvert^2$, and
Lemma \ref{lem:A} gives $\lvert w \rvert^2 = T/9$ for an integer $T$, so
$a = 4T/9$ and $v_2(a) = 2 + v_2(T) \ge 2$: that is $4 \mid a$, with no
integrality of $a$ used. (When $a$ is an integer the same line reads
$9 \mid T$.)

\emph{Backward.} If $4 \mid a$ then $64 \mid S$, so $32 \mid S$, which Lemma
\ref{lem:B} forbids for all-odd coordinates; by \eqref{eq:cong} they are all even,
and with $A = 2A'$ and so on, $S = 4S'$ with $16 \mid S'$ and $A'B' + C'D' = 0$.
The halves satisfy (i): Lemma \ref{lem:C} gives an even number of odd entries
among them, and two is impossible. Suppose exactly two of $A',B',C',D'$ are odd.
The pairs $\{A',B'\}$ and $\{C',D'\}$ are out at once: either would make one side
of $A'B' = -C'D'$ odd and the other even. That leaves the four \emph{mixed}
pairs, one entry from $\{A',B'\}$ and one from $\{C',D'\}$, and they all behave
the same way, for a reason visible in the coefficients of
$S' = A'^2 + 33B'^2 + 3C'^2 + 11D'^2$:
\[
  1 + 3 = 4, \qquad 1 + 11 = 12, \qquad 33 + 3 = 36, \qquad 33 + 11 = 44,
\]
every one of them $\equiv 4 \pmod 8$. An odd square is $\equiv 1 \pmod 8$, so the
two odd terms of $S'$ contribute $4$ modulo $8$. For the two even ones,
$A'B' = -C'D'$ forces them to have the same $2$-adic valuation $k \ge 1$---the odd
entry on each side contributes nothing---so together they contribute
$2^{2k}(\alpha b^2 + \beta d^2)$ with $b,d$ odd and $\{\alpha,\beta\}$ the
complementary pair of coefficients, which is again $\equiv 4 \pmod 8$. At $k = 1$
that is $4 \cdot (8j+4) \equiv 0 \pmod{16}$, and at $k \ge 2$ the factor
$2^{2k}$ already kills it. Either way $S' \equiv 4 \pmod 8$, so $S'$ is $4$ or
$12$ modulo $16$ and never $0$, against $16 \mid S'$.
They satisfy (ii): when all odd, reducing $A'B' = -C'D'$ mod $4$ with every entry
in $\{1,3\}$ leaves an odd number congruent to $3$, which is exactly (ii); when
all even it reduces to $A''+B''+C''+D''$ being even, and $16 \mid S'$ gives
$4 \mid S''$, so Lemma \ref{lem:C} supplies it. Hence $r \in 2\Lat$.
\hfill\qedsymbol

\subsection*{Step 4: every other $\lambda$ imposes nothing new}
Both $r$ and $s$ lie in $V = \Q(\sqrt{33}, i\sqrt3)$, which
Proposition~\ref{prop:rank} says is a field, so $\lambda = s/r$ lies in $V$ too;
being real it lies in $V \cap \R = \Q(\sqrt{33})$, the only real quadratic
subfield. Write $\lambda = \alpha/\beta$ with $\alpha,\beta$ coprime in $\Ok$---which
is possible because $\Ok$ is a principal ideal domain, $\Q(\sqrt{33})$ having class
number one---and put $v = r/\beta$. It lies in $\Lat$: B\'ezout gives
$x\alpha + y\beta = 1$, and then $v = xs + yr \in \Lat$ since $\Lat$ is an
$\Ok$-module, with $\beta v = r$ and $\alpha v = s$. Substituting $r = \beta v$, $s = \alpha v$
into $\lvert u \rvert^2 = 1$ and clearing denominators,
\begin{equation}\label{eq:star}
  \lvert v \rvert^2 \big[\, a(\alpha-\beta)^2 + (4a-1)\alpha\beta \,\big] = a .
\end{equation}

\begin{lemma}\label{lem:val}
For non-zero $v \in \Lat$ and $\pi$ a prime of $\Ok$ above $2$,
$v_\pi(\lvert v \rvert^2) = 2\,v_\pi(v)$, where $v_\pi(v)$ is the largest $k$ with
$v \in \pi^k\Lat$.
\end{lemma}

\begin{proof}
First, $9\lvert x \rvert^2 \in \Ok$ for every $x \in \Lat$---Lemma \ref{lem:A} is
its rational case---because $144\lvert\,\cdot\,\rvert^2$ is a quadratic form with
values in $\Ok$, and both its values and its polarisations on the $30$ unit
vectors are divisible by $16$ there: $30$ checks and $\binom{30}{2} = 435$
(\texttt{step5\_signs.py}). So
$v_\pi\big(144\lvert x\rvert^2\big) \ge 4 = v_\pi(144)$ throughout $\Lat$.

Now write $v = \pi^k v_1$ with $\pi \nmid v_1$. Since $\pi$ is real,
$\lvert v \rvert^2 = \pi^{2k}\lvert v_1 \rvert^2$, so it is enough that
$\lvert v_1 \rvert^2$ be a $\pi$-unit, that is that
$v_\pi\big(144\lvert v_1\rvert^2\big) = 4$ exactly. Replacing $v_1$ by
$v_1 + \pi w$ changes $144\lvert v_1 \rvert^2$ by
$144\pi(v_1\bar w + \bar v_1 w) + 144\pi^2\lvert w \rvert^2$, of valuation at least
$5$ by the previous paragraph, so a valuation of $4$ depends only on the class of
$v_1$ in $\Lat/\pi\Lat \cong \F_2^2$. Its three non-zero classes are represented by the six
points of least norm, $\lvert v \rvert^2 = 1/9$---namely $\pm\tfrac13$ and
$\pm\tfrac16 \pm \tfrac{i\sqrt3}{6}$---and there
$144\lvert v\rvert^2 = 16 = (\pi_1\pi_2)^4$, of valuation exactly $4$.
\end{proof}

\begin{lemma}\label{lem:E}
Let $\pi$ be a prime of $\Ok$ above $2$, and let the denominator of $a$ be
odd---which holds automatically in the only case Corollary~\ref{cor:even} uses,
namely $a$ even, since then the denominator of $a$ in lowest terms is odd; more
generally Theorem~\ref{prop:shells} forces it whenever the shell is non-empty,
since it puts the denominator in $\{1,3,9\}$. If $\pi$ divides none of $v$,
$\alpha$, $\beta$, then $a$ is odd.
\end{lemma}

\begin{proof}
The residue field $\Ok/\pi$ is $\F_2$, so $\alpha \equiv \beta \equiv 1$ there,
hence $\alpha - \beta \equiv 0$ and $\alpha\beta \equiv 1$, and the bracket in
\eqref{eq:star} is $\equiv 4a-1 \equiv 1 \pmod \pi$: writing $a = n/b$ with $b$
odd makes $4n - b$ odd, so $v_\pi(4a-1) = 0$. (With $4 \mid b$ that fails, which
is why the hypothesis is stated.) So
$\pi$ does not divide it, and \eqref{eq:star} gives
$v_\pi(\lvert v \rvert^2) = v_\pi(a) = v_2(a)$. With $\pi \nmid v$, Lemma
\ref{lem:val} makes the left side $0$.
\end{proof}

\begin{corollary}\label{cor:even}
If $4 \nmid a$ then neither prime divides $v$. If moreover $a$ is even then
$\pi_i$ divides exactly one of $r = \beta v$ and $s = \alpha v$, so by
\eqref{eq:colour} exactly one of $\gamma_i(r), \gamma_i(s)$ vanishes and
$\varphi(u) = 0 + c$ with $c \ne 0$: the vector imposes no condition on a
candidate colouring both summands the same way---which is all Step 5 asks
for---for every $\lambda$, rational or not.
\end{corollary}

\begin{proof}
Write $a = n/b$ in lowest terms and clear $b$ in \eqref{eq:star}:
\[
  \lvert v \rvert^2 \big[\, n(\alpha-\beta)^2 + (4n-b)\alpha\beta \,\big] = n ,
\]
a bracket that lies in $\Ok$ whatever the parities, where the bracket of
\eqref{eq:star} need not. By Lemma \ref{lem:val} the factor $\lvert v\rvert^2$ has
valuation $2v_\pi(v)$, so the left side has valuation $2v_\pi(v) + v_\pi(\text{bracket})$;
the bracket lies in $\Ok$, so that second term is at least $0$, and comparing with
the right gives $2v_\pi(v) \le v_2(n)$---an inequality, which is all the proof
uses. If $b$ is odd then
$v_2(n) = v_2(a) \le 1$ because $4 \nmid a$; if $b$ is even then $n$ is odd and
$v_2(n) = 0$. Either way $v_\pi(v) = 0$.

Now let $a$ be even. Lemma \ref{lem:E} puts $\pi$ into one of $v, \alpha, \beta$;
the first is excluded, and $\alpha,\beta$ are coprime, so exactly one of them is
divisible---say $\pi \mid \beta$. Then $\pi \mid r$, while $\alpha \equiv 1$ and
$v \not\equiv 0$ mod $\pi$ leave $\gamma_i(s) = \gamma_i(v) \ne 0$.
\end{proof}

The first half of the corollary settles the whole of Steps 4 and 5 by itself, for
every $\lambda$ at once, and it is worth saying so before the case analysis that
found it. Fix a prime $\pi_i$. Since $\Ok/\pi_i = \F_2$, an $\alpha$ the prime
misses satisfies $\alpha \equiv 1$, so $\alpha v \equiv v$ modulo $\pi_i\Lat$;
hence $\gamma_i(s) = \gamma_i(\alpha v)$ is $0$ if $\pi_i \mid \alpha$ and
$\gamma_i(v)$ if not, and likewise $\gamma_i(r) = \gamma_i(\beta v)$. As
$\alpha,\beta$ are coprime the prime misses at least one of them, leaving three
cases and only three, for $4 \nmid a$:
\begin{itemize}
\item $\pi_i \mid \beta$. Then $\gamma_i(r) = 0$ while
      $\gamma_i(s) = \gamma_i(v) \ne 0$, so $\varphi(u) \ne 0$ for every
      candidate: the vector imposes nothing at all.
\item $\pi_i \mid \alpha$. The same with $r$ and $s$ exchanged.
\item Neither. Then $\gamma_i(r) = \gamma_i(s) = \gamma_i(v)$ and the pair is the
      diagonal $(c,c)$ of Step 2, with $c \ne 0$.
\end{itemize}
So a cross vector with $\lambda \ne -1$ imposes nothing the shell displacements do
not, whatever $\lambda$ is---rational or not, $a$ odd or even. That disposes of
every $\lambda$, at odd and even $a$ alike, and the proof can go straight to the
conclusion. The finer analysis the trichotomy replaced---which $\lambda$ actually
occur at odd $a$---and the sweep that checks both against every rank-$8$ spindle
to $a = 80$ are in the repository, as \texttt{docs/odd-case.tex} and
\texttt{scripts/lambda\_sweep.py}; nothing here depends on either.

\subsection*{Step 5: conclusion}

\begin{lemma}\label{lem:D}
For a shell point $r$, either both primes divide $r$---which happens exactly when
$4 \mid a$---or neither does.
\end{lemma}

\begin{proof}
Both divide if and only if $r \in \pi_1\Lat \cap \pi_2\Lat = 2\Lat$, which by
Step 3 is $4 \mid a$. Suppose exactly one does, say $r = \pi_1 w$. Then
$\lvert w \rvert^2 = a/\pi_1^2 = a(29 - 5\sqrt{33})/8$, so in the chart $w$ has
\[
  S = 522a, \qquad 2(AB+CD) = -90a .
\]
A non-empty shell puts $a$ in $\tfrac19\Z$ by Lemma \ref{lem:A}, and $9$ divides
both $522$ and $90$, so these are integers; it also makes $v_2(a) \ge 0$, so the
two cases below exhaust $4 \nmid a$. If $a$ is odd then $S \equiv 2 \pmod 4$, and
$S \equiv [A]+[B]+3[C]+3[D]$ takes the value $2$ only on the patterns
$\mathrm{ooee}$ and $\mathrm{eeoo}$, both of which violate (i). If
$v_2(a) = 1$, write $a = 2a'$ with $a'$ odd; then $S \equiv 4 \pmod 8$,
which all-odd coordinates cannot give since they force $S \equiv 0 \pmod 8$, so all
coordinates are even and $AB + CD = 4(A'B' + C'D')$, while $AB + CD = -45a = -90a'$
is $2 \pmod 4$. The case $r = \pi_2 w$ is the image of this one under the
Galois conjugation $\sigma$ of $F$, which swaps $\pi_1$ and $\pi_2$ and changes
only the sign of $90a$; every congruence above is a statement modulo $4$ or $8$
about $S = 522a$ and $AB+CD = \mp45a$, and none of them sees that sign.
\end{proof}

If $4 \mid a$ \emph{and the shell is non-empty}---this is the only place the
second hypothesis is used, and the only half of the theorem that needs it---Step 3
gives a shell point $r = 2w$, and then $(1-\theta_a)w$ lies in $\Lam_a$ with length
$\tfrac12$, so Proposition~\ref{prop:half} rules out both groups of order $4$ at
once. Concretely, the shell displacement $u = (1-\theta_a)r$ has both colours zero,
so $\varphi(u) = 0$ for all $24$ Klein candidates and none survives. If $4 \nmid a$, Lemma
\ref{lem:D} says no shell point is divisible by either prime, so both colourings
give it a non-zero colour: the shell displacements impose only diagonal pairs
$(c,c)$ with $c \ne 0$, uniformly over the shell and for either choice of
colouring. Fix one of the two and use it on both summands, which is what makes
the pairs diagonal at all. By Step 4 nothing else constrains, and two automorphisms
differing at every non-zero point satisfy $\psi_1(c) \ne \psi_2(c)$---the relative
permutation must be a $3$-cycle, and two of the six are---so at least four of the
twenty-four candidates survive, and a colouring exists. This proves Theorem \ref{thm:law}.
\hfill\qedsymbol

\section{The distance spectrum, and the candidates it leaves}\label{sec:seven}

The law leaves $4 \mid a$ with a non-empty shell as the candidates
(Figure~\ref{fig:law}). Which $a$ have
one is decided by the form \eqref{eq:norm} and not by a ball, and has a closed
answer, because $\Lat$ is $\tfrac13\Ok_V$ and occupancy is a norm question.

Call a prime $p \ne 3, 11$ \emph{obstructed} when it splits in $\Q(\sqrt{33})$ and
stays inert in $\Q(\sqrt{-3})$. For odd $p$ that is the pair of symbols
\[
  \Big(\tfrac{33}{p}\Big) = 1
  \qquad\text{and}\qquad
  \Big(\tfrac{-3}{p}\Big) = -1,
\]
and neither need be computed: since $33 \equiv 1 \pmod 4$, quadratic
reciprocity gives $\big(\tfrac{33}{p}\big) = \big(\tfrac{p}{33}\big)$, and
$\big(\tfrac{-3}{p}\big) = 1$ exactly when $p \equiv 1 \pmod 3$. Both are
therefore fixed by $p$ modulo $33$, and running the $20$ classes coprime to $33$
leaves five:
\[
  p \text{ is obstructed}
  \iff
  p \equiv 2,\ 8,\ 17,\ 29 \text{ or } 32 \pmod{33}.
\]
The prime $2$ is obstructed too, straight from the definition---$33 \equiv 1
\pmod 8$ splits it, $-3 \equiv 5 \pmod 8$ keeps it inert---and it lands in the
class $2 \pmod{33}$ with the others, so the residue form covers it as well and
nothing below makes it a separate case. Below $200$ the obstructed primes are
$2\mc 17\mc 29\mc 41\mc 83\mc 101\mc 107\mc 131\mc 149\mc 167\mc 173$ and $197$; the residue form is
checked against the symbols for every odd prime below $20000$.

\begin{remark}[The five classes are a genus]\label{rem:narrow}
The five classes are not an artefact of running twenty. Since $h(F) = 1$ and the
fundamental unit has norm $+1$, the narrow class number of $F = \Q(\sqrt{33})$ is
$h^+(F) = 2$, so $F$ has exactly one quadratic extension unramified at every
finite prime. It can be named without any computation: $\operatorname{disc}(F) =
33 = (-3)\cdot(-11)$ is a product of two prime discriminants, so the narrow genus
field of $F$ is $\Q(\sqrt{-3},\sqrt{-11})$, and since $\sqrt{-3}\sqrt{-11} =
-\sqrt{33}$ that field is $V$. So

\begin{center}
$V$ is the narrow Hilbert class field of $\Q(\sqrt{33})$.
\end{center}

\noindent
This gives $\mathfrak d_{V/F} = (1)$ a second time, by a route independent of the
discriminant comparison in Proposition~\ref{prop:order}; the two agree, and
either may be taken as the source.

What it buys is that the local behaviour used above and below is class field
theory rather than case work: a prime of $F$ splits in $V$ exactly when it has a
\emph{totally positive} generator. The three cases the paper needs fall out at
sight. The primes above $2$ are generated by $(\pm5\pm\sqrt{33})/2$, of norm
$-2$, whose embeddings $+5.372$ and $-0.372$ have opposite signs; no unit can
repair that, since the units are $\pm\varepsilon^k$ with
$\Nm(\varepsilon) = +1$---so $\pi_1$ and $\pi_2$ are inert, which is what
Section~\ref{sec:arith} found by hand. $\mathfrak p_3 = (6+\sqrt{33})$ embeds as
$+11.745$ and $+0.255$, is totally positive, and splits.
$\mathfrak p_{11} = (11+2\sqrt{33})$ embeds as $+22.489$ and $-0.489$ and does
not.

Over $\Q$ the same statement is genus theory of discriminant $33$. The genus
characters are $(\tfrac{\cdot}{3})$ and $(\tfrac{\cdot}{11})$, there are two
classes of binary quadratic forms of that discriminant and one in each genus, and
$p$ is obstructed exactly when both characters are $-1$: \emph{$p$ lies in the
non-principal genus of discriminant $33$.} That is the same five classes, and it
explains why five. Half of the twenty classes coprime to $33$ fail the first
character, and half of those fail the second.
\end{remark}

\begin{theorem}\label{prop:shells}
Let $a = n/b$ be a positive rational in lowest terms. Then $\Lat$ has a point at
squared radius $a$ if and only if $b \in \{1,3,9\}$ and $v_p(n)$ is even at every
obstructed $p$.
\end{theorem}

\begin{proof}
By Proposition~\ref{prop:order}, $a$ is a squared radius of $\Lat$ if and only if
$9a = \Nm_{V/F}(w)$ for some $w \in \Ok_V$. Write $\varepsilon = 23 - 4\sqrt{33}$, a generator of the totally positive units
of $F$---the unit group is $\langle -1\rangle \times \langle\varepsilon\rangle$ and
$-1$ is negative at both real places, so those units are exactly the powers of
$\varepsilon$. Three further facts are used, and it is worth saying where each
comes from.

$V/F$ is cyclic of degree $2$ and unramified at every finite prime: that is the
by-product of Proposition~\ref{prop:order}'s discriminant comparison, not a new
assumption.

$h(F) = 1$, in two lines. The Minkowski bound for $F = \Q(\sqrt{33})$ is
$\tfrac12\sqrt{33} < 2.88$, so every ideal class contains an integral ideal of
norm at most $2$, that is the unit ideal or a prime above $2$. Both of those are
principal, since $2 = -\pi_1\pi_2$ with $\pi_{1,2} = (5 \pm \sqrt{33})/2$ in
$\Ok_F$. So the class group is trivial.

$\varepsilon = 23 - 4\sqrt{33}$ generates the totally positive units of $F$---it
is the inverse of the fundamental unit $23 + 4\sqrt{33} \approx 45.98$ under the
usual convention, and generates the same group---and it is totally
positive: $23^2 - 33\cdot4^2 = 1$, and both conjugates $23 \mp 4\sqrt{33}$ are
positive. Since $-1$ is negative at both real places, the totally positive units
are exactly the powers of $\varepsilon$.

$h(V) = 1$, by Minkowski again, and the computation is small enough to give in
full. With $n = 4$, $r_2 = 2$ and $\lvert\operatorname{disc}(V)\rvert = 1089$ the
bound is
\[
  \Bigl(\frac4\pi\Bigr)^{2}\frac{4!}{4^4}\sqrt{1089}
  = \frac{16}{\pi^2}\cdot\frac{3}{32}\cdot 33 < 5.02 ,
\]
so every ideal class contains an integral ideal of norm at most $5$. Which primes
of $\Ok_V$ have norm that small? Above $2$: each $\pi_i$ is inert in $V$, so the
prime above it has norm $\Nm(\pi_i)^2 = 4$, and there is no prime of norm $2$.
Above $3$: the prime $\mathfrak p_3$ of $F$ ramifies, and splits in $V$ into two
primes of norm $3$. Above $5$: $33 \equiv 3 \pmod 5$ is a non-residue, so $5$ is
inert in $F$ and every prime above it has norm at least $25$. Hence the only
integral ideals of norm at most $5$ are the unit ideal, the two primes of norm $3$
and the two of norm $4$, and it is enough to generate those four.

The ones above $2$ are $\pi_1\Ok_V$ and $\pi_2\Ok_V$, principal already in
$\Ok_F$. For those above $3$, take
\[
  w = \tfrac12\bigl(1 + i(2\sqrt3 + \sqrt{11})\bigr) \in \Ok_V,
\]
which lies in $3\Lat$ because $w/3$ is the chart point $(2,0,4,2)$ of
\eqref{eq:chart}---the point $(6,0,12,6)$ verifies only $w \in \Lat$---and has
\[
  \Nm_{V/F}(w) = \lvert w \rvert^2
  = \tfrac14\bigl(1 + (2\sqrt3+\sqrt{11})^2\bigr)
  = \tfrac14\bigl(24 + 4\sqrt{33}\bigr) = 6 + \sqrt{33},
\]
so $\Nm_{V/\Q}(w) = 36 - 33 = 3$ and $(w)$ is one of the two; its conjugate under
$\sqrt{33} \mapsto -\sqrt{33}$ is the other. So the class group is trivial.

(\textsc{Pari} agrees, and \texttt{scripts/moser\_is\_ov.py} runs the check, but
the argument above does not need it.)

Which primes of $F$ are inert in $V$? Exactly those above an obstructed $p$,
together with the one above $11$. That is the content of the two symbols: $p$
splits in $F$ and stays inert in $\Q(\sqrt{-3})$ precisely when the prime above it
does not split in $V$.

\emph{Necessity.} Let $9a = \Nm(w)$ with $w \in \Ok_V$. Then $9a$ lies in
$\Ok_F \cap \Q = \Z$, so $b \mid 9$. At a prime $P$ of $F$ inert in $V$,
$v_P(\Nm(w)) = 2v_{P'}(w)$ is even. For $P$ above an obstructed $p$ we have
$v_P = v_p$ on rationals, so $v_p(9a)$ is even; as no obstructed $p$ equals $3$,
$v_p(a)$ is even too. (Above $11$ the prime ramifies in $F$, so $v_P = 2v_{11}$ is
even for every rational and the condition is empty.)

\emph{Sufficiency.} Suppose $b \mid 9$ and $v_p(a)$ is even at every obstructed
$p$. Then $9a$ is a positive rational, hence totally positive in $F$, and has even
valuation at every prime of $F$ inert in $V$. At the two infinite places $V$ is
totally imaginary, so the local norms there are the positive reals, which $9a$
meets; at a finite prime that splits there is no condition, and at an inert one the
condition is even valuation, so $9a$ is a local norm everywhere and Hasse's norm
theorem for the cyclic extension $V/F$ gives $9a = \Nm(x)$ for some
$x \in V^\times$. We record that because it is the natural first move, but the
argument does not use it: what is wanted is an algebraic \emph{integer}, and the
route to one goes through the ideal directly, with no appeal to Hasse at any
point.

Build the ideal, then. Even valuation at the inert
primes is exactly what lets $(9a)\Ok_F$ be written as $\Nm_{V/F}(A)$ for an ideal
$A$ of $\Ok_V$; since $h(V) = 1$, $A = (w)$, and $\Nm(w)$ generates $(9a)$, so
$\Nm(w) = 9a\,u$ for a unit $u$ of $\Ok_F$. Both $9a$ and $\Nm(w) = \lvert w\rvert^2$
are totally positive, so $u$ is. The totally positive units of $F$ are the powers
of $\varepsilon$---the unit group is $\langle -1\rangle \times \langle
\varepsilon\rangle$ and $-1$ is negative at both places---and $\varepsilon = \Nm_{V/F}(\eta)$ for
\[
  \eta = i\bigl(2\sqrt3 - \sqrt{11}\bigr) \in \Ok_V,
\]
which is purely imaginary, so $\bar\eta = -\eta$ and
\[
  \Nm_{V/F}(\eta) = \eta\bar\eta = \bigl(2\sqrt3 - \sqrt{11}\bigr)^2
  = 12 + 11 - 4\sqrt{33} = 23 - 4\sqrt{33} = \varepsilon,
\]
while $\Nm_{V/\Q}(\eta) = 23^2 - 16\cdot 33 = 1$, so $\eta$ is a unit. It lies in
$\Ok_V = 3\Lat$ because $\eta/3 = i(2\sqrt3 - \sqrt{11})/3$ is the chart point
$(0,0,8,-4)$. Replacing $w$
by $w\eta^{-k}$ with $u = \varepsilon^{k}$ gives $\Nm(w) = 9a$ exactly.
\end{proof}

The ingredients are all small and all checkable: $h(F) = h(V) = 1$,
$\operatorname{disc}(V) = 1089$, the signature $(0,2)$, and
$\Nm_{V/F}(\eta) = \varepsilon$, which the accompanying code verifies along with
$3\Lat$ being the maximal order.
Where the parity comes from is set out below, after the density.

The denominator contributes nothing beyond being admissible, which is consistent:
$3$ is not obstructed, and the reason is local rather than a matter of
ramification---$-11$ is a square in $\Q_3$, so $\sqrt{-3}$ already lies in
$\Q_3(\sqrt{33})$ and the extension splits there, leaving no condition. At $11$ the
extension is inert, but $11$ ramifies in $\Q(\sqrt{33})$, so a rational has even
valuation at the prime above it and the condition is empty again. The restriction $b \in \{1,3,9\}$ is Lemma~\ref{lem:A}'s corollary from
Section~\ref{sec:proof}, and with it the criterion covers \emph{every} rational
radius, not merely the integers. The three rational spindles found below are all
$n = 64$ or $256$, powers of two with even
exponent, and all three have non-empty shells.

Put beside Theorem~\ref{thm:law} this describes the occupied-shell candidates off
the factorisation of $a$, with no ball and no search. It does not say which
candidates admit a 5-chromatic join: Section~\ref{sec:stronger}
kills eleven of the sixteen below $100$; $a = 64$ is the one it cannot reach,
and Lemma~\ref{lem:chain} settles that one by construction instead.

\begin{corollary}\label{cor:closed}
Let $a \ge 1/2$, normalised as in the introduction, and write $a = n/b$ in
lowest terms. Then $\theta_a$ is a candidate---the law admits it and its shell is
non-empty---if and only if
\[
  b \in \{1,3,9\},
  \qquad
  v_p(n) \text{ is even at every obstructed } p,
  \qquad
  v_2(n) \ge 2 .
\]
\end{corollary}

Only the last clause is new: $2$ being obstructed, the middle clause already asks
for $v_2$ even, and $4 \mid a$ adds $v_2 \ge 2$. That single line accounts for the
nine empty multiples of $4$ below $100$ without any computation---eight have
$v_2$ odd, and the ninth is $68 = 2^2\cdot17$, where the obstructed prime $17$
appears once---and it was checked against the form for every integer $a \le 400$
without exception.

It also settles how common candidates are. Writing $\mathcal C$ for the integer
candidates, each condition restricts at one prime and distinct primes are
independent, so for any finite set $P$ of obstructed primes containing $2$,
\[
  \overline d(\mathcal C) \;\le\; \frac16 \prod_{\substack{p \in P \\ p \ne 2}}
  \frac{p}{p+1},
\]
the $\tfrac16$ being $2$'s own share. Five of the twenty classes mod $33$ are
obstructed, so those primes have Dirichlet density $\tfrac14$, $\sum 1/p$
diverges over them, and the product tends to $0$: the integer candidates have
natural density zero. The measured densities at $N = 10^2$, $10^3$ and $10^4$ and the Dirichlet
density of the obstructed primes are checked in \texttt{audit\_numbers.py}.

The rule that produced every row is deterministic and is stated here so that the
table can be read, even though the constructions themselves belong to the
companion paper. Each half is a \emph{whole ball}: all points of $\Lat$ within a
radius, at the smallest depth in the thirty unit vectors that reaches the shell
at $\sqrt a$, with the two halves glued at the origin and the second turned by
$\theta_a$. No vertex was chosen, no search was run over subsets, and no
parameter was tuned---which is why the vertex counts are large, and why they are
a property of that rule rather than of the rotation. Two rows were afterwards
minimised, and the archive's \texttt{scripts/ball\_params.py} recovers each
row's depth and radius from the shipped coordinates rather than from a build
log, so the table can be checked against the graphs instead of trusted.

\begin{center}
\begin{tabular}{lllrrr}
\hline
rotation & $\cos$ & field & depth & radius & vertices ($\chi \ge 5$) \\
\hline
$\theta_4$ & $7/8$ & $\Q(\sqrt3,\sqrt5,\sqrt{11})$ & --- & --- & $509$, minimised \\
$\theta_{16}$ & $31/32$ & $\Q(\sqrt3,\sqrt7,\sqrt{11})$ & $4$ & $4.0000$ & $29113$ \\
$\theta_{28}$ & $55/56$ & $\Q(\sqrt3,\sqrt{11},\sqrt{37})$ & $6$ & $5.3113$ & $156577$ \\
$\theta_{36}$ & $71/72$ & $\Q(\sqrt3,\sqrt{11},\sqrt{13})$ & $6$ & $6.0000$ & $175825$ \\
$\theta_{64/3}$ & $125/128$ & $\Q(\sqrt3,\sqrt{11},\sqrt{23})$ & $7$ & $4.6383$ & $187225$ \\
$\theta_{64/9}$ & $119/128$ & $\Q(\sqrt3,\sqrt{11},\sqrt{247})$ & $3$ & $2.6817$ & $5809$ \\
$\theta_{256/9}$ & $503/512$ & $\Q(\sqrt3,\sqrt{11},\sqrt{1015})$ & $6$ & $5.3506$ & $158257$ \\
\hline
\end{tabular}
\end{center}

The first is Parts' record and was known; the four of \cite{Voronov} are the
other previously known rotations, and $\theta_4$ is common to both lists. Of the
six new ones, the three rational spindles follow from Parts' two published
monochromatic-pair gadgets by the chaining below: his $367$-vertex gadget has
terminals $8/3$ apart and gives $a = 64n^2/9$, whose cases $n = 1, 2$ are
$\theta_{64/9}$ and $\theta_{256/9}$; his $421$-vertex one has them $8/\sqrt3$
apart and gives $a = 64n^2/3$, whose case $n = 1$ is $\theta_{64/3}$. So
$\theta_{16}$, $\theta_{28}$ and $\theta_{36}$ are the three for which we found
no prior construction. Every row carries a DRAT certificate that \texttt{drat-trim}
checks, shipped in the archive at the commit pinned in \emph{Code and data}
below, so the verdicts here stand on the archive and not on any unpublished
work: eight proofs, all
reporting \texttt{s VERIFIED}, re-checkable in about ten minutes by
\texttt{scripts/check\_drat.sh}.

The two lemmas the introduction leans on are short enough to keep.

\begin{lemma}[chaining a monochromatic pair]\label{lem:chain}
Let $M \subset \Lat$ be finite with two marked points $0$ and $t \ne 0$ such that
every proper $4$-colouring of $M$ gives them the same colour, and let $n \ge 1$ be
such that $a = n^2\lvert t\rvert^2$ is rational and exceeds $\tfrac14$, so that
$\theta_a$ is one of the rotations of Section~\ref{sec:intro}. Put
$M_n = \bigcup_{j=0}^{n-1}(M + jt)$. Then $M_n \subset \Lat$, every proper
$4$-colouring of $M_n$ gives $0$ and $nt$ the same colour, and
$\lvert M_n \rvert \le n(\lvert M \rvert - 1) + 1$. Consequently, writing
$a = n^2\lvert t \rvert^2$,
\[
  M_n \cup \theta_a(M_n)
\]
is not $4$-colourable, on at most $2n(\lvert M \rvert - 1) + 1$ vertices.
\end{lemma}

\begin{proof}
$\Lat$ is a group and $t \in \Lat$, so every translate $M + jt$ lies in $\Lat$ and
so does the union; $a > \tfrac14$ makes $\sqrt{4a-1}$ real, so $\theta_a$ exists,
and $a \in \Q$ makes it a spindle. In a proper $4$-colouring the copy $M + jt$
forces $jt$ and
$(j+1)t$ to agree, so $0$ and $nt$ agree; the copies overlap in at least the
shared marked point, which gives the count. The same argument inside
$\theta_a(M_n)$ forces $\theta_a(0) = 0$ and $\theta_a(nt)$ to agree, and the two
copies share the origin, so $nt$ and $\theta_a(nt)$ both carry the colour of the
origin. But $\lvert nt \rvert^2 = a$ puts $nt$ on the shell, so
$\lvert nt - \theta_a(nt) \rvert = 1$ and they are adjacent.
\end{proof}

\begin{lemma}\label{lem:subjoin}
Let $G = L \cup \theta(S)$ be a finite non-$4$-colourable join sharing the origin.
Then some induced subgraph of $G$ containing the origin is a join of the same
form and has chromatic number exactly $5$.
\end{lemma}

\begin{proof}
Among the induced subgraphs of $G$ that contain the origin and are not
$4$-colourable---a non-empty family, since $G$ is one---take one, $H$, with
fewest vertices. Then $\chi(H) \ge 5$. For the other direction pick any
$v \ne 0$ of $H$; by minimality $H - v$ has a proper $4$-colouring, and giving
$v$ a fifth colour extends it to $H$, so $\chi(H) \le 5$. And $H$ is a join of the
same form: putting
\[
  L_H = L \cap V(H), \qquad S_H = \theta^{-1}\bigl(V(H) \cap \theta(S)\bigr),
\]
we have $V(H) = L_H \cup \theta(S_H)$ with $L_H \subseteq L$ and
$S_H \subseteq S$, and $0 \in V(H)$ by construction.
\end{proof}

Lemma~\ref{lem:chain} applied to Parts' $367$-vertex gadget, whose two terminals
lie $8/3$ apart, gives a join at $a = 64n^2/9$ for every $n \ge 1$; its premise
is the gadget's monochromatic-pair property, which is itself one of the eight
certificates. His inventory lists a second monochromatic pair, at $8/\sqrt3$ on
$421$ vertices, and the same lemma sends it to $a = 64n^2/3$.

\section{A stronger filter, and what survives it}\label{sec:stronger}

The criterion of Section~\ref{sec:ball} uses group homomorphisms, which are
exactly the \emph{linear} proper colourings of $\Lat/m\Lat$. Every proper
colouring of that finite quotient pulls back to the whole lattice, linear or not,
so using all of them is strictly stronger and still needs no ball. A periodic
colouring sees a cross pair only through its residues, so all it can see is which
\emph{residue pairs} occur, and $(\alpha,\beta)$ occurs exactly when the coset
$(\alpha - \theta\beta) + m\Lam$ contains a unit vector of the joint lattice.

One degeneracy must be excluded: if $(0,0)$ occurs---equivalently if $\Lam$ holds
a vector of length $1/m$---both halves give the centre's coset the same colour
and the test cannot kill. At $m = 2$ that is precisely the half-unit criterion of
Section~\ref{sec:closed}. The moduli tried are $4$, $6$, $8$ and $12$, and the
rule followed is to report the smallest of those at which $(0,0)$ does not occur
\emph{and} a kill is found; a rotation degenerate at every one of them is set
aside as undecided rather than as a survivor. Non-degeneracy alone does not
suffice: $\theta_{48}$ is non-degenerate at $m = 5$, where $72$ pairs occur, and
survives there.

Applied to the integer spindles to $a = 100$, it kills $\theta_a$ for $a = 12$,
$20$, $44$, $48$, $52$, $60$, $76$, $80$, $84$, $92$ and $100$---nine at $m = 4$
over $4^4 = 256$ residue classes, and $a = 48$, $80$, degenerate at $m = 4$ and
$m = 6$, at $m = 8$ over $4096$. Each kill ships as a witness naming its modulus,
with both colourings and every occurring residue pair, re-checked to be proper,
to agree at the shared origin and to disagree on every occurring pair. That is
the whole infinite union, not a sample of it.

That leaves $\theta_4$, $\theta_{16}$, $\theta_{28}$, $\theta_{36}$---and
$\theta_{64}$, where $(0,0)$ occurs at every modulus tried, so the filter never
speaks. It never could have: Lemma~\ref{lem:chain} exhibits a join at $a = 64$,
and no colouring argument excludes what exists. The degeneracy is forced rather
than unlucky, since $\Lat \cap \Q = \tfrac13\Z$ makes $(8/m)(1-\theta_{64})$ a
vector of length $1/m$ in $\Lam_{64}$ whenever $8/m \in \tfrac13\Z$, which holds
for every $m \mid 24$; the four moduli tried all divide $24$.

The empty-shell multiples of $4$ lie outside this tally, since the law is silent
there, and are settled separately. There are nine below $100$. Six die: $a = 8$,
$32$, $40$, $56$, $68$ and $96$ each carry a homomorphism $4$-colouring---exactly
$96$ apiece, $24$ Klein and $72$ cyclic, counted as in
Section~\ref{sec:arith}---so no subset of the joint lattice is 5-chromatic there.
Three survive: $a = 24$, $72$ and $88$ admit none, and the reason is
Proposition~\ref{prop:half} rather than anything in this section---each carries
exactly twelve vectors of length $\tfrac12$, one of $\Lam_{24}$'s written out in
Section~\ref{sec:closed}, so no homomorphism onto a group of order $4$ can avoid
the unit vectors whatever the shell does. An earlier version of the accompanying
code declined to run the periodic filter at an empty shell, on the ground that
the question did not arise. It does arise, and these three are what the refusal
was hiding. Each of the nine ships a witness that
\texttt{scripts/check\_emptyshell.py} validates from its JSON alone.

\section{What is not claimed}\label{sec:notclaimed}

The criterion is a one-sided kill. A survivor is a candidate; only an explicit
5-chromatic union makes it a construction, which is why Section \ref{sec:seven}
exhibits graphs rather than citing survivals.

The law is about \emph{this} lattice, and the same criterion on Haugland's
lattice $L_1$, generated by the $42$nd roots of unity, is the evidence:
$\theta_2$ dies on the Moser lattice and survives there, as does $\theta_7$.
Those computations are \texttt{hept\_rank.py}, \texttt{hept\_classify.py} and
\texttt{hept\_halfunit.py}. No attempt is made to transfer the law.

The rank-$8$ hypothesis excludes rotations that map $\Lat$'s span into
itself---$\theta_3$ is $\omega_3$, $\theta_1$ is $\omega_1$---where the cross
decomposition does not arise. Proposition~\ref{prop:rank} says exactly which
those are, $4a-1 = 3k^2$ or $11k^2$, so the hypothesis is recognisable at sight,
and proved rather than assumed.

Proposition \ref{prop:half} is unconditional for both groups of order $4$. An
earlier version was not, and the remark in Section~\ref{sec:closed} records why:
its three exceptions were a bug in the sweep, not a feature of the lattice.

Theorem~\ref{prop:shells} is proved, through
Proposition~\ref{prop:order}: $\Lat$ is $\tfrac13\Ok_V$, so occupancy is a
question about relative norms, which the arithmetic of $\Ok_V$ answers---by
$h(V) = 1$, even valuation at the inert primes and the unit $\eta$, and not, as
the proof says, by Hasse. An earlier version of this paper stated it as verified
and not proved, having missed that $3\Lat$ is the maximal order. Theorem~\ref{thm:law} is likewise proved, and the
constructions of Section~\ref{sec:seven} are exhibited rather than predicted, so
nothing there rests on the criterion.

The choice of $\Lat$ is a choice. Taking $\lvert j\rvert \le J$ instead of
$\lvert j\rvert \le 2$ in the generating family gives, for each $J$, a rank-$4$
lattice, and these nest: $30 \subset 42 \subset 54 \subset \cdots$. The index over
ours is $9^{J-2}$, checked for $J \le 5$, and the unit vectors are exactly the
$6(2J+1)$ generators, checked for $J \le 3$. Neither is proved here for all $J$;
counting generators does not rule out a further unit vector, and nothing here
depends on the general statement.
Heule's $553$-vertex large part lives in the $42$-vector member. Everything here
is about the $30$-vector one, which the record and every construction in
Section~\ref{sec:seven} use; the law is not claimed for the finer members, and
the criterion would have to be recomputed on each.

Nothing here reaches the rotations of \cite{Voronov} that are not spindles; we
checked rather than assumed it. Their four are $\rho u$ for a unit vector $u$
of $\Lat$; the whole family $\{\rho u\}$ is $30$ rotations, of which six---the
sixth roots of unity---fix $\Lat$ and so give back $\theta_4$'s own joint lattice,
and exactly two are spindles, $\rho$ and $-\rho$. Every one of the $30$ has a
rank-$8$ joint lattice with $126$, $108$ or $90$ unit vectors. Run against all
$30$, the criterion of Section \ref{sec:ball} kills \emph{none}, and the periodic
filter of Section \ref{sec:stronger} none either, degenerating at $m = 4$ and
$m = 8$ alike. So both tools are blind outside the spindle family: against the
integer spindles the law removes three $a$ in four and the periodic filter eleven
of the sixteen survivors; here they remove nothing. Whatever decides that family
is not in this paper.

No claim is made about the record: $509$ is unbeaten---de~Grey was still
looking for a way past it in September $2026$ \cite{Thread18}---and nothing here
brings it closer. What is settled is an obstruction---no spindle of rank $8$ with
$4 \nmid a$ can be used at all, the rank-$4$ exceptions of
Proposition~\ref{prop:rank} being untouched---together with a proved criterion for
occupancy,
eleven further exclusions and seven constructions. What is not settled is the
classification those would have to add up to: rotations
outside the spindle family are untouched.

\section*{Code and data}

Every computation here is carried out by a script in the archive,
\url{https://github.com/calinnedelcu/joining-rotations-artifacts}, whose
\texttt{README.md} maps each claim to the script that reproduces it and says what
the archive settles and what it does not. Its code and data stand at commit
\texttt{78b7579a74373cdb9fd1391ffcd1b278a27c26c9}.
\texttt{audit\_numbers.py} re-derives the paper's numbers from the definitions in
three minutes on the standard library alone, sharing no code with the rest---bar
the nine empty-shell verdicts, which \texttt{empty\_shell\_status.py} derives and
it checks against;
\texttt{check\_emptyshell.py} and \texttt{check\_periodic.py} check the twenty
witnesses against their JSON;
\texttt{check\_drat.sh} verifies the eight DRAT certificates---every construction
of the companion paper, the gadget Lemma~\ref{lem:chain} rests on, and Parts'
record---in ten minutes, \texttt{drat-trim} deliberately not vendored; and
\texttt{verify\_all.sh} runs
the whole suite in parallel, about an hour on six cores; its last full run passed
$35$ of $35$, and the eleven entries added since pass on their own. \texttt{PROVENANCE.md}
names the four files taken from \cite{Parts}, \texttt{LICENSES.md} the licence split: MIT for the
code, CC BY 4.0 for the paper, certificates and data. A companion paper on the
seven constructions and the $\theta_4$ interface bound is in preparation; nothing
above depends on it, and its computations are already in the archive. Archived at
Zenodo, \doi{10.5281/zenodo.22855504}---the \emph{concept} DOI, which follows the
newest deposit; the commit above fixes the snapshot.

\section*{Acknowledgements}

The homomorphism $4$-colouring of the Moser ring, which this paper's criterion
redirects at the joint lattice, is due to Philip Gibbs in the Polymath16 project;
the $\omega_t$ notation and the algebraic framework are Polymath16's. The
classification of the Moser lattice's two geometric $4$-colourings is due to
\'Akos D\'ucz.

% The bibliography ends a few lines short of holding the address block, which
% would otherwise sit alone on a twenty-fifth page.
\enlargethispage{3\baselineskip}
\begin{thebibliography}{99}
\setlength{\itemsep}{0pt}\setlength{\parsep}{0pt}

\bibitem{deGrey}
A.~D.~N.~J. de~Grey,
\emph{The chromatic number of the plane is at least 5},
Geombinatorics \textbf{28} (2018), 18--31; arXiv:1804.02385.

\bibitem{Ducz}
\'A. D\'ucz,
\emph{A note on geometric colorings of the Moser lattice},
arXiv:2606.12325 (2026).

\bibitem{Engel}
P. Engel, O. Hammond-Lee, Y. Su, D. Varga, P. Zs\'amboki,
\emph{Diverse beam search to find densest-known planar unit distance graphs},
arXiv:2406.15317.

\bibitem{Ruhland}
H. Ruhland,
\emph{Families of lattices with an unbounded number of unit vectors},
arXiv:2410.16172.

\bibitem{Haugland}
J.~K. Haugland,
\emph{A Moser-spindle-free 5-chromatic unit distance graph on 2131 vertices in
the plane},
arXiv:2608.04542 (2026).

\bibitem{Heule}
M.~J.~H. Heule,
\emph{Computing small unit-distance graphs with chromatic number 5},
Geombinatorics \textbf{28} (2018), 32--50; arXiv:1805.12181.

\bibitem{GibbsRing}
P. Gibbs, \emph{Elements of the Moser spindle ring}, preprint, May 2018,
the author's own deposit;
\url{https://www.researchgate.net/publication/324950826_Elements_of_the_Moser_Spindle_Ring}.

\bibitem{Hubai}
T. Hubai,
\emph{Comment on Polymath16, fourth thread: applying the probabilistic method},
Polymath16, 2018;
\url{https://dustingmixon.wordpress.com/2018/05/05/polymath16-fourth-thread-applying-the-probabilistic-method/}.

\bibitem{Parts}
J. Parts,
\emph{Graph minimization, focusing on the example of 5-chromatic unit-distance
graphs in the plane},
Geombinatorics \textbf{29} (2020), 137--166; arXiv:2010.12665.

\bibitem{Matolcsi}
M. Matolcsi, I.~Z. Ruzsa, D. Varga, P. Zs\'amboki,
\emph{The fractional chromatic number of the plane is at least 4},
arXiv:2311.10069 (2023).

\bibitem{DuczVarga}
\'A. D\'ucz, D. Varga,
\emph{A unit-distance graph in the plane with independence ratio below 1/4},
arXiv:2606.28157 (2026).

\bibitem{Speyer}
D. Speyer, comment of 3 May 2018 on \emph{Polymath16, third thread: is
6-chromatic within reach?}, Polymath16, 2018;
\url{https://dustingmixon.wordpress.com/2018/05/01/polymath16-third-thread-is-6-chromatic-within-reach/}.

\bibitem{Polymath}
Polymath16,
\emph{Algebraic formulation of the Hadwiger--Nelson problem},
Polymath wiki, 2018.

\bibitem{Thread18}
A.~D.~N.~J. de~Grey,
comments of 13 September 2026 on \emph{Polymath16, eighteenth thread: back with
a new conjecture}, a thread opened 17 January 2026;
\url{https://dustingmixon.wordpress.com/2026/01/17/polymath16-eighteenth-thread-back-with-a-new-conjecture/}.

\bibitem{Voronov}
V. Voronov, A. Neopryatnaya, E. Dergachev,
\emph{Constructing 5-chromatic unit distance graphs embedded in the Euclidean
plane and two-dimensional spheres},
Discrete Math. \textbf{345} (2022), 113106; arXiv:2106.11824.

\end{thebibliography}

\end{document}
